\documentclass[11pt,a4paper]{article}
\usepackage[T1]{fontenc}
\usepackage[utf8]{inputenc}
\IfFileExists{lmodern.sty}{\usepackage{lmodern}}{\usepackage{mathptmx}}
\usepackage[french]{babel}
\usepackage[margin=2.3cm]{geometry}
\usepackage{amsmath,amssymb,mathtools,bm}
\usepackage{graphicx,booktabs,array,multirow,longtable}
\usepackage{xcolor}
\usepackage{tikz}
\usetikzlibrary{arrows.meta,positioning,shapes.geometric,fit,calc}
\usepackage{tcolorbox}
\usepackage{enumitem}
\usepackage{caption}
\usepackage{float}
\usepackage[hidelinks]{hyperref}
\captionsetup{font=small,labelfont=bf}
\usepackage{siunitx}
\sisetup{output-decimal-marker={,},group-separator={\,}}
\usepackage{algorithm}
\usepackage{algpseudocode}
\floatname{algorithm}{Algorithme}
\definecolor{bleu}{HTML}{2a78d6}
\definecolor{orange}{HTML}{eb6834}
\definecolor{vert}{HTML}{1baf7a}
\definecolor{fond}{HTML}{f4f7fb}
\newtcolorbox{encadre}[1][]{colback=fond,colframe=bleu!60,boxrule=0.5pt,arc=2pt,left=6pt,right=6pt,top=4pt,bottom=4pt,#1}
\newtcolorbox{alerte}[1][]{colback=orange!6,colframe=orange!70,boxrule=0.5pt,arc=2pt,left=6pt,right=6pt,top=4pt,bottom=4pt,#1}
\newcommand{\dd}{\mathrm{d}}
\newcommand{\ii}{\mathrm{i}}
\newcommand{\hatt}[1]{\widehat{#1}}
\newcommand{\cs}{\texttt{chausspec}}
\renewcommand{\epsilon}{\varepsilon}

\title{\textbf{Influence de la texture de surface sur la réponse\\ des chaussées aéronautiques}\\[6pt]
\large Cadre mécaniste à deux échelles (contact pneumatique--texture et calcul multicouche
viscoélastique sous charge roulante)\\ et comparaison avec les résultats du TFE}
\author{Lucas \textsc{David}\\ \small LTDS -- ENTPE, thèse en partenariat avec le STAC}
\date{25 septembre 2026 --- rapport de travail, version 2.0}

\begin{document}
\sloppy
\maketitle

\begin{abstract}
Le TFE a modélisé la chaussée de la plateforme PEP sous un bogie d'A340 avec une surface
\emph{lisse} : la pression de contact y varie seulement à l'échelle du pneumatique. Une piste
réelle est texturée (macrotexture de 1 à \SI{50}{mm}) et souvent rainurée. Le pneumatique n'y
touche que les sommets de la texture, et la pression locale dépasse largement la pression
nominale. Ce rapport construit, vérifie et exploite un cadre mécaniste qui propage cet effet
jusqu'aux déformations de la chaussée. Il comprend trois étapes. (1) La génération de surfaces
contrôlées : rainurage normalisé FAA, et macrotextures aléatoires de profondeur moyenne de profil
(MPD) imposée selon l'ISO 13473-1. (2) Le calcul du contact entre la bande de roulement et la
texture, par une méthode d'éléments de frontière spectrale vérifiée contre deux solutions de
référence. (3) Un couplage deux échelles avec le calcul semi-analytique multicouche
viscoélastique de \cs, qui respecte le fait que la texture est fixe dans la chaussée alors que
l'enveloppe de pression roule. Chaque hypothèse de la chaîne est contrôlée par un calcul
indépendant, et le rapport précise pour chaque résultat le degré de confiance qu'on peut lui
accorder.
La version 2 ajoute un second niveau de modélisation : le frottement au freinage (calcul 3D
complet), l'entaille géométrique des rainures et la microstructure granulats--mortier, ces deux
derniers par un modèle local éléments finis chargé par les grandeurs du calcul 3D au point
critique.
Conclusion principale : la texture amplifie fortement les déformations des premiers millimètres
de la couche de roulement ; le frottement, l'entaille des rainures et la microstructure
accentuent encore cet effet. Cet effet est confiné aux 10 à \SI{20}{mm} supérieurs. Les résultats du
TFE en profondeur (base de couche liée, cisaillement maximal vers \SI{100}{mm}) ne sont pas
modifiés.
\end{abstract}

\tableofcontents
\newpage

%=============================================================================================
\section{Question posée et réponse courte}
%=============================================================================================
\textbf{Question.} La texture de la surface (rugosité, rainurage) rend-elle la sollicitation de
la chaussée plus sévère que le cas lisse étudié dans le TFE, et à quelle profondeur ?

\noindent\textbf{Niveau 1} (enrobé homogène, surface plane, sans frottement ; \S\ref{sec:cadre} à \S\ref{sec:resultats}) :
\begin{encadre}
\textbf{Réponse courte.}
\begin{enumerate}[nosep]
\item \textbf{Oui près de la surface.} Sous l'empreinte hétérogène du TFE, la déformation
      principale majeure maximale est multipliée par 1,6 à 2,0 à \SI{1}{mm}
      de profondeur, par 1,5 à 1,8 à \SI{2}{mm} et par 1,3 à 1,5 à \SI{5}{mm},
      selon la texture (rainurage FAA, macrotextures de MPD 0,6 à \SI{1.5}{mm}). Les
      cisaillements proches de la surface sont amplifiés davantage encore.
\item \textbf{Non en profondeur.} Le facteur tombe à 1,05--1,23 à \SI{10}{mm} et à
      1,00--1,01 à \SI{20}{mm} ; l'effet est nul à \SI{40}{mm} (1,00--1,00). La déformation en base de couche liée et le cisaillement
      maximal vers \SI{100}{mm}, qui sont les deux grandeurs clés du TFE, ne sont pas modifiés.
\item La texture crée donc une \textbf{seconde zone critique, superficielle} (0 à
      \SI{10}{mm}), distincte de celle du TFE. C'est la zone d'amorçage de la fissuration par le
      haut et de l'arrachement.
\item L'effet croît avec la MPD, mais \textbf{la MPD ne suffit pas} : à MPD égale, la pente des
      aspérités, l'asymétrie de la texture et surtout la raideur du caoutchouc changent le
      résultat.
\item \textbf{Confiance} : les tendances et les profondeurs d'influence sont solides ; la valeur
      absolue aux premiers millimètres est incertaine d'un facteur de l'ordre de 1,5, à cause de
      paramètres mal connus (\S\ref{sec:verif}).
\end{enumerate}
\end{encadre}
\begin{encadre}[colframe=orange!70]
\textbf{Ce que change le niveau 2 (frottement, entaille des rainures, granulats et mortier).}
\begin{enumerate}[nosep]
\item \textbf{Les conclusions du niveau 1 tiennent} : amplification près de la surface, effet nul
      au-delà de 20 à \SI{40}{mm}, résultats du TFE en profondeur inchangés.
\item \textbf{Mais le niveau 1 sous-estimait la sévérité superficielle}, surtout pour les
      rainures. L'entaille multiplie encore $\epsilon_1$ par 2,1 à \SI{1}{mm}
      (arêtes vives ; 1,6 avec des arêtes arrondies), avec deux points chauds :
      le bord supérieur des plateaux et le fond de la rainure. Avec un freinage courant ($\mu=0{,}3$), une surface rainurée
      atteint 3,7 à
      4,0 fois la déformation
      de la surface lisse à \SI{1}{mm}, contre 1,7 au niveau 1.
\item \textbf{Le frottement} augmente $\epsilon_1$ superficielle de 46\,\%
      (surface lisse) à 90\,\% (MPD \SI{1.5}{mm})
      pour $\mu=0{,}6$ : il accentue l'effet de la texture.
\item \textbf{Dans le mortier} entre granulats, la déformation locale sous une texture vaut
      2,7 à 3,7 fois la déformation moyenne de l'enrobé.
      Cette concentration existe aussi, et plus forte, sous une surface lisse : c'est une propriété
      du matériau. L'écart \emph{relatif} entre texture et surface lisse est donc atténué dans le
      mortier (facteur 0,65 à 0,77),
      alors que la déformation \emph{absolue} y est bien plus forte.
\item \textbf{Confiance} : les sens de variation sont sûrs ; les facteurs du niveau 2 sont
      semi-quantitatifs (facteur 1,5 à 2), car ils dépendent de paramètres à mesurer : distance
      critique, arrondi des arêtes, contraste granulat/mortier (\S\ref{sec:n2_confiance}).
\end{enumerate}
\end{encadre}


\paragraph{Organisation du rapport.} La section~\ref{sec:cadre} fixe les échelles et les
hypothèses. La section~\ref{sec:surfaces} explique en détail la construction des surfaces, et la
section~\ref{sec:contact} le calcul du contact. La section~\ref{sec:couplage} décrit le passage du
contact à la chaussée, étape par étape. La section~\ref{sec:verif} rassemble toutes les
vérifications et en tire un degré de confiance. La section~\ref{sec:resultats} présente les
résultats et la comparaison avec le TFE. La section~\ref{sec:niveau2} (nouvelle dans cette
version) lève les hypothèses de frottement nul, de surface plane et d'enrobé homogène, et corrige
les résultats en conséquence. La section~\ref{sec:discussion} discute les limites et la
suite. Les annexes donnent les paramètres numériques et la procédure de reproduction.

%=============================================================================================
\section{Cadre, échelles et hypothèses}\label{sec:cadre}
%=============================================================================================
\subsection{Trois échelles}
\begin{center}\small
\begin{tabular}{p{2.6cm}p{2.3cm}p{4.3cm}p{5.2cm}}
\toprule
Échelle & Longueur & Objets & Traitement dans ce rapport\\
\midrule
Structure / pneu & 0,3--\SI{5}{m} & bogie, empreintes, couches & calcul multicouche VE sous charge roulante (\cs), identique au TFE\\
Macrotexture & 1--\SI{50}{mm} & granulats, rainures & surfaces générées, contact caoutchouc / surface rigide, réponse locale\\
Microtexture & $<\SI{0.5}{mm}$ & aspérités des granulats & ignorée (rôle dans l'adhérence, pas dans la réponse structurelle)\\
\bottomrule
\end{tabular}
\end{center}
Le rapport d'échelle entre la longueur de contact ($L\approx\SI{0.56}{m}$) et la période de la
texture ($\Lambda\lesssim\SI{50}{mm}$) vaut au moins 10. C'est ce qui justifie une approche à deux
échelles.

\subsection{Chaîne de calcul}
\begin{figure}[H]
\centering
\begin{tikzpicture}[node distance=0.55cm and 0.6cm, >=Stealth,
  box/.style={draw=bleu!70, fill=fond, rounded corners=2pt, align=center, font=\small, minimum height=1.0cm, text width=3.6cm},
  res/.style={draw=orange!80, fill=orange!6, rounded corners=2pt, align=center, font=\small, minimum height=1.0cm, text width=3.6cm}]
\node[box] (surf) {\textbf{1. Surface} $h(x,y)$\\rainures FAA ou texture aléatoire (MPD cible)};
\node[box, right=of surf] (cont) {\textbf{2. Contact} bande de roulement / surface rigide, pour 8 pressions nominales};
\node[box, right=of cont] (lib) {\textbf{3. Bibliothèque}\\$m(x,y;\bar p)=p/\bar p$};
\node[box, below=of surf] (nom) {\textbf{4. Réponse nominale}\\bogie A340, KVG, $V$ = 0,66 m/s (TFE, \cs)};
\node[box, right=of nom] (hist) {\textbf{5. Histoire de} $\delta p(x,s)$\\en chaque point fixe de la chaussée};
\node[box, right=of hist] (her) {\textbf{6. Transformation héréditaire}\\$\delta p_{\mathrm{eq}}=E_s\!\int J\,\dd_s\delta p$};
\node[box, below=of her] (loc) {\textbf{7. Réponse locale}\\massif homogène, grille 0,25 mm};
\node[res, left=of loc] (sum) {\textbf{8. Somme}\\$\bm\epsilon=\bm\epsilon_0+\delta\bm\epsilon$\\$\epsilon_1$, $\epsilon_{xz}$, $\epsilon_{yz}$};
\node[res, left=of sum] (cmp) {\textbf{9. Comparaison}\\lisse (TFE) / texturé, en fonction de $z$};
\draw[->] (surf) -- (cont); \draw[->] (cont) -- (lib); \draw[->] (lib) -- (hist);
\draw[->] (nom) -- (hist); \draw[->] (hist) -- (her); \draw[->] (her) -- (loc); \draw[->] (loc) -- (sum);
\draw[->] (nom.south) -- ++(0,-0.27) -| (sum.north); \draw[->] (sum) -- (cmp);
\end{tikzpicture}
\caption{Chaîne de calcul. Les étapes 1 à 3 sont à l'échelle de la texture, l'étape 4 à l'échelle
du pneu ; les étapes 5 à 7 couplent les deux échelles.}\label{fig:chaine}
\end{figure}

\subsection{Hypothèses}
\begin{enumerate}[label=(H\arabic*),leftmargin=*,nosep]
\item \textbf{Chaussée rigide pour le contact.} $E_{\text{enrobé}}\approx\SI{15000}{MPa}$ contre
      $E_{\text{caoutchouc}}\approx\SI{10}{MPa}$. La géométrie du contact ne dépend donc pas de la
      déformation de la chaussée ; l'erreur relative est de l'ordre de $10^{-3}$.
\item \textbf{Bande de roulement} élastique linéaire, de module $E_r=\SI{10}{MPa}$,
      $\nu_r=0{,}49$ et épaisseur $t_r=\SI{10}{mm}$, collée à une ceinture indéformable. Valeurs
      \emph{supposées}, à confirmer pour un pneu 1400$\times$530R23 ; sensibilité étudiée à 5 et
      \SI{20}{MPa}.
\item \textbf{Contact normal sans frottement}\,\textsuperscript{$\dagger$}, quasi-statique (le temps de passage vaut
      \SI{0.8}{s} et le caoutchouc est supposé répondre instantanément).
\item \textbf{Séparation d'échelles.} Sur une cellule de texture, la pression nominale $p_0$ est
      localement uniforme ; le contact s'y établit comme sous $\bar p=p_0$.
\item \textbf{Texture fixe dans la chaussée.} Seule l'enveloppe $p_0(x-Vt,y)$ roule ; la
      modulation de contact suit la position fixe de la texture.
\item \textbf{Chaussée à surface plane}\,\textsuperscript{$\dagger$}. La texture n'intervient que par la modulation de la
      pression. La géométrie (entaille des rainures, creux) n'est pas représentée.
\item \textbf{Couche de roulement homogène}\,\textsuperscript{$\dagger$} (enrobé équivalent), viscoélastique linéaire (KVG du
      TFE), à coefficient de Poisson constant (0,35).
\end{enumerate}
(H6) et (H7) sont les deux limites structurelles du niveau 1. \textsuperscript{$\dagger$}\,Les
hypothèses (H3), (H6) et (H7) sont levées par le niveau 2 (\S\ref{sec:niveau2}), qui corrige les
résultats du niveau 1.

%=============================================================================================
\section{Construction des surfaces}\label{sec:surfaces}
%=============================================================================================
\subsection{Rainurage transversal normalisé}
\paragraph{Géométrie.} La FAA (AC 150/5320-12C) prescrit des rainures sciées de
\SI{6}{mm} de largeur ($w=\SI{6.35}{mm}$, soit 1/4") et \SI{6}{mm} de profondeur, au pas de
\SI{38}{mm} ($\lambda=\SI{38.1}{mm}$, soit 1,5"). Elles sont perpendiculaires à l'axe de piste,
donc orientées selon $y$ dans nos axes ($x$ = sens de roulement). Le profil ne dépend que de $x$ :
\begin{equation}
h(x)=\begin{cases}-d_g & \text{si } |u|<w/2\\ 0 & \text{sinon}\end{cases},\qquad
u=\operatorname{mod}\!\left(x-\phi+\tfrac\lambda2,\lambda\right)-\tfrac\lambda2 ,
\end{equation}
où $\phi$ fixe la position des rainures dans la chaussée. On prend $\phi=0$ : comme on balaie
toute la zone de contact, tous les décalages entre rainures et empreinte sont représentés.
\paragraph{Arrondi des arêtes.} Une arête vive crée une pression singulière au bord des plats
(\S\ref{sec:contact_ref}). Pour mesurer la sensibilité à cette idéalisation, on remplace l'arête
par un quart de cercle de rayon $r$ : pour une distance $s\in[0,r]$ à la paroi, côté plat,
$h=-\bigl(r-\sqrt{r^2-(r-s)^2}\bigr)$. On étudie $r=0$ (arête vive) et $r=\SI{1}{mm}$ (arête
usée ou épaufrée).
\paragraph{Discrétisation.} Un pas est discrétisé en 152 points (\SI{0.25}{mm}). La résolution
est étudiée jusqu'à \SI{0.0625}{mm} (\S\ref{sec:resolution}).

\begin{figure}[H]\centering\includegraphics[width=0.78\linewidth]{figures/m_rainures.pdf}
\caption{Haut : profil d'un pas de rainurage. Bas : pression de contact correspondante sous
\SI{1.65}{MPa} nominal (bande de roulement de \SI{10}{mm}). Avec arête vive, la surpression au
bord est singulière : son maximum n'est fixé que par la résolution (figure calculée au pas de \SI{0.0625}{mm} ; \SI{7.2}{MPa} au pas de \SI{0.25}{mm} de la campagne).}\label{fig:rainures}
\end{figure}

\subsection{Macrotexture aléatoire : principe}
Une macrotexture de chaussée est un champ aléatoire de hauteurs. Sa description minimale, standard
en mécanique du contact, est sa \emph{densité spectrale de puissance} (DSP) $C(\bm q)$ : la part
de la variance des hauteurs portée par chaque nombre d'onde $\bm q$. On construit donc des
surfaces à DSP imposée, puis on règle leur amplitude pour obtenir une MPD cible. Ce procédé a
trois avantages :
\begin{itemize}[nosep]
\item la surface est \textbf{contrôlée} : on sait exactement quelles longueurs d'onde elle
      contient ;
\item elle est \textbf{périodique}, ce qui convient à la FFT et permet de la répéter sous tout le
      pneu ;
\item elle est \textbf{reproductible} : une graine de tirage fixe la surface, et l'étude de
      plusieurs graines mesure la dispersion d'une réalisation à l'autre.
\end{itemize}

\subsection{Macrotexture aléatoire : algorithme de génération}
\begin{algorithm}[H]
\caption{Génération d'une macrotexture de MPD cible (\texttt{contact.random\_texture})}
\begin{algorithmic}[1]
\State Grille périodique $N\times N$ de côté $\Lambda_c$ ; ici $\Lambda_c=\SI{128}{mm}$ et
       $N=512$, soit un pas $\Delta=\SI{0.25}{mm}$.
\State Nombres d'onde de la FFT : $q_x,q_y=2\pi\,\texttt{fftfreq}(N,\Delta)$ ; $q=\sqrt{q_x^2+q_y^2}$.
\State DSP isotrope auto-affine avec plateau :
       $C(q)=1$ pour $q<q_r$ et $C(q)=(q/q_r)^{-2(1+H)}$ sinon, avec $q_r=2\pi/\SI{20}{mm}$ et
       $H=0{,}8$ ; $C=0$ hors de $[q_0,q_1]=[2\pi/\SI{50}{mm},\,2\pi/\SI{1}{mm}]$ et $C(0)=0$.
\State Phases aléatoires indépendantes, uniformes : $\theta(\bm q)\sim\mathcal U[0,2\pi)$ (graine fixée).
\State $g=\operatorname{Re}\,\mathrm{FFT}^{-1}\!\left[\sqrt{C(q)}\,e^{\ii\theta}\right]$, puis
       centrage et réduction : $g\leftarrow(g-\bar g)/\sigma_g$.
\State Asymétrie : $h=g$ (gaussienne) ; $h=e^{\beta g}$ (positive) ; $h=-e^{\beta g}$ (négative),
       avec $\beta=0{,}4$ ; puis centrage.
\State Calcul de la MPD de $h$ (algorithme~\ref{alg:mpd}), selon $x$ et selon $y$ ; moyenne des deux.
\State Mise à l'échelle : $h\leftarrow h\times\text{MPD}_{\text{cible}}/\text{MPD}(h)$ (la MPD est
       homogène de degré 1).
\end{algorithmic}
\end{algorithm}
\paragraph{Justification des choix.}
\begin{itemize}[nosep]
\item \textbf{Bornes 1--\SI{50}{mm}} : c'est la définition normalisée de la macrotexture
      (ISO 13473-1 : 0,5--\SI{50}{mm}). La borne basse est portée à \SI{1}{mm} pour être résolue
      par au moins 4 points à \SI{0.25}{mm} ; la part de variance perdue est négligeable, car la
      DSP décroît en $q^{-3,6}$.
\item \textbf{Plateau au-delà de \SI{20}{mm}} : les longueurs d'onde supérieures à la taille des
      plus gros granulats ne portent pas plus d'énergie. Le plateau évite qu'elles dominent
      artificiellement.
\item \textbf{$H=0{,}8$} : valeur typique des surfaces de chaussée rapportée dans la littérature
      du contact pneu/chaussée (Persson) ; c'est un paramètre à caler sur des relevés réels.
\item \textbf{Asymétrie} : les enrobés présentent plutôt une texture « négative » (plateaux et
      creux, $S_{sk}<0$), les enduits et gravillonnages une texture « positive » (pics,
      $S_{sk}>0$). La transformation exponentielle, monotone, conserve la corrélation spatiale et
      règle $S_{sk}\approx\pm1{,}3$.
\end{itemize}

\begin{figure}[H]\centering\includegraphics[width=0.75\linewidth]{figures/m_psd.pdf}
\caption{DSP imposée (trait) et DSP de la surface générée (moyenne radiale, cercles). Les pointillés
marquent les bornes de 50, 20 et \SI{1}{mm}.}\label{fig:psd}
\end{figure}

\subsection{Calcul de la MPD (ISO 13473-1)}
\begin{algorithm}[H]
\caption{Profondeur moyenne de profil (\texttt{contact.mpd\_iso13473})}\label{alg:mpd}
\begin{algorithmic}[1]
\For{chaque ligne de la surface (profil selon $x$ ; puis de même selon $y$)}
  \For{chaque segment de base de \SI{100}{mm}}
    \State retirer la droite de régression (redressement du segment) ;
    \State couper le segment en deux moitiés et relever le pic de chaque moitié ;
    \State profondeur du segment $=\tfrac12(\text{pic}_1+\text{pic}_2)-$ niveau moyen du segment.
  \EndFor
\EndFor
\State MPD = moyenne de toutes les profondeurs de segment.
\end{algorithmic}
\end{algorithm}
La profondeur de texture équivalente, comparable à l'essai volumétrique à la tache de sable,
vaut $\text{ETD}=0{,}2+0{,}8\,\text{MPD}$ (mm). Les MPD étudiées (0,6, 1,0 et \SI{1.5}{mm})
correspondent à des ETD de 0,68, 1,0 et \SI{1.4}{mm}. Pour mémoire, l'OACI recommande une
profondeur de texture d'au moins \SI{1.0}{mm} pour une piste neuve, et la FAA \SI{1.14}{mm}.

\begin{figure}[H]\centering\includegraphics[width=0.8\linewidth]{figures/m_mpd.pdf}
\caption{Calcul de la MPD sur un segment de \SI{100}{mm} de la surface gaussienne (MPD cible
\SI{1.0}{mm} ; ce segment vaut \SI{0.87}{mm}, la moyenne sur tous les segments vaut exactement 1,0).}
\end{figure}
\begin{figure}[H]\centering\includegraphics[width=0.9\linewidth]{figures/m_textures.pdf}
\caption{Les trois familles de texture à MPD égale (\SI{1.0}{mm}) : cartes des hauteurs et
histogrammes. À MPD identique, la répartition des sommets est très différente.}\label{fig:textures}
\end{figure}

\paragraph{Pourquoi la MPD ne suffit pas.} La MPD ne mesure que la hauteur des pics par rapport au
niveau moyen. Pour le contact, ce qui compte est la \emph{pente} des aspérités, mesurée par
$m_2=\langle|\nabla h|^2\rangle$ : à charge donnée, l'aire de contact est inversement
proportionnelle à $\sqrt{m_2}$ (\S\ref{sec:persson}). Or, à MPD égale, $m_2$ varie du simple au
triple entre nos textures (tableau~\ref{tab:surfaces}). Un indicateur mécanique de texture devra
donc combiner hauteur et pente, par exemple MPD et $m_2$, ou exploiter la DSP complète.

\begin{table}[H]\centering\small
\caption{Caractéristiques des surfaces et du contact sous la pression nominale du TFE (\SI{1.65}{MPa}).
$m_2$ : pente quadratique moyenne ; aire : fraction de la surface en contact ; $p_{99}$ : 99\textsuperscript{e}
centile de la pression locale dans le contact (pour les rainures : maximum, qui dépend de la
résolution à cause de la singularité d'arête).}\label{tab:surfaces}
\begin{tabular}{lccccccc}
\toprule
Surface & MPD & ETD & $S_q$ & $S_{sk}$ & $m_2$ & aire & $p_{99}$\\
 & (mm) & (mm) & (mm) & & & (\%) & (MPa)\\
\midrule
rainures FAA, arêtes vives & --- & --- & --- & --- & --- & 84 & 7,2 (max)\\
rainures FAA, arêtes $r$ = 1 mm & --- & --- & --- & --- & --- & 82 & 4,3 (max)\\
aléatoire, MPD 0,6 mm & 0,60 & 0,68 & 0,32 & $-$0,03 & 0,050 & 76 & 5,5\\
aléatoire, MPD 1,0 mm & 1,00 & 1,00 & 0,53 & $-$0,03 & 0,140 & 55 & 8,2\\
aléatoire, MPD 1,5 mm & 1,50 & 1,40 & 0,80 & $-$0,03 & 0,315 & 40 & 11,7\\
négative, MPD 1,0 mm & 1,00 & 1,00 & 0,74 & $-$1,28 & 0,289 & 52 & 7,2\\
positive, MPD 1,0 mm & 1,00 & 1,00 & 0,41 & 1,28 & 0,089 & 70 & 9,5\\
MPD 1,0 ; $E_r$ = 5 MPa & 1,00 & 1,00 & 0,53 & $-$0,03 & 0,140 & 83 & 4,8\\
MPD 1,0 ; $E_r$ = 20 MPa & 1,00 & 1,00 & 0,53 & $-$0,03 & 0,140 & 31 & 15,2\\
\bottomrule
\end{tabular}
\end{table}


%=============================================================================================
\section{Contact bande de roulement / texture}\label{sec:contact}
%=============================================================================================
\subsection{Formulation}
Soit $h(x,y)$ la hauteur de la surface rigide, positive vers le haut, et $u(x,y)\ge0$ le
déplacement normal de la surface du caoutchouc, positif en compression. La surface du caoutchouc
se trouve en $d+u$, où $d$ est le rapprochement rigide. L'écartement vaut $g=u-h+d$ et le problème
s'écrit
\begin{equation}
p\ge0,\qquad g\ge0,\qquad p\,g=0\quad\text{(conditions de Signorini)},\qquad
\langle p\rangle=\bar p\quad\text{(équilibre)} ,
\end{equation}
avec $u=\mathcal C*p$ (élasticité linéaire du caoutchouc). C'est la condition d'optimalité de la
minimisation de l'énergie complémentaire
$\Pi(p)=\tfrac12\int p\,(\mathcal C*p)-\int p\,h$ sur le cône $p\ge0$ sous la contrainte
$\langle p\rangle=\bar p$. $\Pi$ est strictement convexe : la solution existe et elle est unique.

\subsection{Souplesse de la bande de roulement}
Dans l'espace de Fourier, $\hatt u(\bm q)=C(q)\,\hatt p(\bm q)$. $C(q)$ est calculée par le
\emph{même noyau multicouche} que \cs : une couche de caoutchouc d'épaisseur $t_r$ collée sur un
substratum rigide, soumise à une pression unitaire. On a les deux limites
\begin{equation}
C(q)\xrightarrow[qt_r\gg1]{}\frac{2(1-\nu_r^2)}{E_rq}\ \ \text{(massif semi-infini)},\qquad
C(q)\xrightarrow[qt_r\ll1]{}\frac{t_r(1+\nu_r)(1-2\nu_r)}{E_r(1-\nu_r)}\ \ \text{(compression confinée)} .
\end{equation}
Les longueurs d'onde courtes de la texture ($\lambda\ll2\pi t_r\approx\SI{60}{mm}$) voient un
massif semi-infini ; les plus longues voient une couche mince quasi incompressible, donc raide.
Le caoutchouc « enveloppe » alors moins les ondulations longues.

\subsection{Résolution numérique}
\begin{algorithm}[H]
\caption{Contact par gradient conjugué projeté (Polonsky \& Keer, 1999 ; \texttt{contact.solve\_contact})}
\begin{algorithmic}[1]
\State $p\leftarrow\bar p$ partout ; $\delta\leftarrow0$ ; $G_{\text{old}}\leftarrow1$
\Repeat
  \State $u=\mathrm{FFT}^{-1}[C\,\mathrm{FFT}(p)]$ ; $g=u-h$ ; $I_c=\{p>0\}$ ; $g\leftarrow g-\operatorname{moy}_{I_c}(g)$
  \State $G=\sum_{I_c}g^2$ ; direction $t=g+\delta\,(G/G_{\text{old}})\,t$ sur $I_c$, $t=0$ ailleurs ; $G_{\text{old}}\leftarrow G$
  \State $r=\mathcal C*t$ ; $r\leftarrow r-\operatorname{moy}_{I_c}(r)$ ; pas $\tau=\sum_{I_c}g\,t/\sum_{I_c}r\,t$
  \State $p\leftarrow\max(p-\tau t,0)$ \Comment{projection sur $p\ge0$}
  \State points hors contact en interpénétration ($p=0$, $g<0$) : $p\leftarrow p-\tau g$, $\delta\leftarrow0$ ; sinon $\delta\leftarrow1$
  \State $p\leftarrow p\,\bar p/\langle p\rangle$ \Comment{équilibre}
\Until{$\sum|p-p_{\text{old}}|/\sum p<10^{-7}$}
\end{algorithmic}
\end{algorithm}
Pour une cellule de $512^2$ points, le calcul demande 100 à 700 itérations, soit 5 à
\SI{20}{s}.

\subsection{Solutions de référence}\label{sec:contact_ref}
\paragraph{Poinçons plans périodiques (rainures).} Tant que le caoutchouc ne touche pas le fond
des rainures, les plats (largeur $2a=\lambda-w$) sont en contact total. Pour un massif
semi-infini, c'est le problème classique des poinçons plans rigides périodiques :
\begin{equation}
\frac{p(x)}{\bar p}=\frac{\cos(\pi x/\lambda)}{\sqrt{\sin^2(\pi a/\lambda)-\sin^2(\pi x/\lambda)}}\quad(|x|<a),\qquad
a_n=\tfrac12\bigl[P_n(c)+P_{n-1}(c)\bigr],\ c=\cos(2\pi a/\lambda),
\label{eq:punch}
\end{equation}
où les $a_n$ sont les coefficients de Fourier ($p/\bar p=a_0+2\sum a_n\cos(2\pi nx/\lambda)$) et
$P_n$ les polynômes de Legendre. Cette distribution est indépendante du module et singulière en
$1/\sqrt{\cdot}$ aux arêtes. Le caoutchouc s'enfonce d'environ \SI{0.8}{mm} dans les rainures,
profondes de \SI{6.35}{mm} : l'hypothèse de contact limité aux plats est vérifiée.
\paragraph{Théorie de Persson (surfaces aléatoires).}\label{sec:persson} Pour un massif
semi-infini, $A/A_0=\operatorname{erf}\!\bigl(\sqrt2\,\bar p/(E_r^*\sqrt{m_2})\bigr)$. On sait
qu'elle sous-estime de 20 à 30\,\% l'aire aux faibles charges.

\subsection{Bibliothèque de contact}
La pression nominale varie sous le pneu, de 0 à \SI{3.4}{MPa} pour l'empreinte hétérogène, et le
contact sur texture aléatoire est \emph{non linéaire} : l'aire augmente avec $\bar p$. On résout
donc le contact d'une même cellule pour 8 pressions nominales (0,15 ; 0,3 ; 0,6 ; 1,0 ; 1,65 ;
2,3 ; 3,0 ; \SI{4.0}{MPa}) et on stocke les modulations $m(x,y;\bar p)=p/\bar p$. En un point de
pression nominale $p_0$, la modulation est interpolée linéairement en $\log p_0$ entre deux
niveaux. Pour les rainures, le contact est total sur les plats et linéaire : $m(x)$ ne dépend pas
de $\bar p$.
\begin{figure}[H]\centering\includegraphics[width=\linewidth]{figures/m_bibliotheque.pdf}
\caption{Bibliothèque de contact de la texture gaussienne (MPD \SI{1.0}{mm}) : l'aire de contact
passe de 13\,\% à 90\,\% entre 0,3 et \SI{4.0}{MPa}.}\label{fig:biblio}
\end{figure}

%=============================================================================================
\section{Du contact à la chaussée : couplage deux échelles}\label{sec:couplage}
%=============================================================================================
\subsection{Décomposition}
La pression réelle à l'instant $s$, au point fixe $\bm x=(x,y)$ de la chaussée, vaut
\begin{equation}
p(\bm x,s)=\underbrace{p_0(x-Vs,y)}_{\text{enveloppe roulante (TFE)}}\;m\bigl(\bm x;p_0(x-Vs,y)\bigr)
= p_0+\delta p,\qquad \delta p=p_0\,(m-1) .
\end{equation}
La modulation $m$ est prise dans la bibliothèque à la position \emph{fixe} $\bm x$, la cellule de
texture étant répétée périodiquement sur toute la chaussée. Par linéarité de la chaussée,
$\bm\epsilon=\bm\epsilon_0+\delta\bm\epsilon$ : la réponse nominale $\bm\epsilon_0$ est celle du
TFE, calculée par \cs, et $\delta\bm\epsilon$ est la contribution de la texture.

\subsection{Réponse nominale}
Bogie d'A340 complet (4 roues de \SI{370}{kN}), structure PEP, KVG du Tableau 7 du TFE,
$V=\SI{0.66}{m/s}$, régime permanent, domaine de $\SI{6}{m}\times\SI{3}{m}$ et grille de
$2048\times1024$ (pas de \SI{2.9}{mm}), aux profondeurs 1, 2, 5, 10, 20, 40 et \SI{80}{mm}.
Les champs sont interpolés sur la grille fine. Ce calcul est celui de la validation V5 de la
notice \cs, qui reproduit le modèle COMSOL du TFE.

\subsection{Localisation : la texture ne « voit » que la couche de surface}
Une composante de $\delta p$ de longueur d'onde $\lambda$ s'amortit en profondeur comme
$(1+qz)e^{-qz}$, avec $q=2\pi/\lambda$ ; sa \emph{profondeur d'influence} vaut
$\lambda/2\pi\le\SI{8}{mm}$. C'est très petit devant l'épaisseur de la couche de roulement
(\SI{320}{mm}) : $\delta p$ charge localement un \textbf{massif homogène semi-infini}. Sa réponse
se calcule sur une grille fine, de pas \SI{0.25}{mm}, couvrant l'empreinte de la roue arrière
droite plus une marge de \SI{2}{cm} ($2720\times2480$ points), par la solution fermée du massif
homogène dans l'espace de Fourier (\texttt{texture.local\_response}). Ce point est vérifié contre
la structure multicouche complète (T6, \S\ref{sec:verif}).

\subsection{Viscoélasticité : méthode héréditaire}\label{sec:hered}
\paragraph{Pourquoi un traitement particulier.} $\delta p$ n'est pas une charge roulante : la
texture est fixe, seule l'enveloppe se déplace. De plus, pour une texture aléatoire, la
\emph{forme} de $\delta p$ change pendant le passage, puisque l'aire de contact suit $p_0$.
\paragraph{Principe.} Pour un massif homogène à coefficient de Poisson constant chargé en effort,
le principe de correspondance donne
\begin{equation}
\delta\bm\epsilon(\bm x,t)=\int_{-\infty}^{t}J(t-s)\,\frac{\partial}{\partial s}\,\delta\bm\epsilon^{E=1}(\bm x,s)\,\dd s,
\qquad \delta\bm\epsilon^{E=1}(\cdot,s)=\mathcal G_1[\delta p(\cdot,s)],
\end{equation}
où $J$ est la complaisance de fluage et $\mathcal G_1$ l'opérateur élastique de module unité.
$\mathcal G_1$ est linéaire en espace et ne dépend pas du temps : il \textbf{commute} avec
l'intégrale d'hérédité. D'où
\begin{equation}
\boxed{\;\delta\bm\epsilon(\bm x,0)=\mathcal G_{E_s}\bigl[\delta p_{\text{eq}}\bigr],\qquad
\delta p_{\text{eq}}(\bm x)=E_s\int_{-\infty}^{0}J(-s)\,\frac{\partial\,\delta p(\bm x,s)}{\partial s}\,\dd s\;}
\label{eq:hered}
\end{equation}
L'intégrale d'hérédité s'applique donc \emph{point par point à la pression}, puis on fait un seul
calcul élastique. Cette réduction est exacte dans le cadre (H1)--(H7), y compris pour le contact
non linéaire et pour la mémoire laissée par la roue avant du bogie au même endroit de la chaussée.
\paragraph{Intégration.} Pour le KVG,
$J(t)=1/E_0+\sum_i(1-e^{-t/\tau_i})/E_i$. On suit l'histoire de $\delta p$ en chaque point depuis
l'arrivée de la roue avant (environ \SI{4}{s} avant l'instant d'observation), par pas de 5 à
\SI{10}{mm} de déplacement de l'enveloppe. On utilise des variables internes récurrentes
$h_i$, avec une intégration exacte pour une variation linéaire de $\delta p$ sur chaque pas. Les
cellules de temps caractéristique très court ($\tau<0{,}02\,\Delta s$) sont traitées comme
instantanées. L'erreur de pas est inférieure à $2\cdot10^{-4}$ pour l'empreinte hétérogène et à
1\,\% pour les empreintes en créneau.

\begin{alerte}
\textbf{Méthode abandonnée en cours de travail.} La première version (bilan du 24/09) utilisait
une formule multiplicative : $\delta\bm\epsilon=\mathcal G_{E_s}[\delta p(t=0)]\cdot E_s\,C(X,y)$,
avec une souplesse effective $C=Q/p_0$. Elle suppose $\delta p(\bm x,s)$ proportionnel à
$p_0(x-Vs)$ en chaque point. C'est exact pour le rainurage (contact total, modulation
fixe) : les deux méthodes y coïncident à 0,2\,\% (T5). Pour une texture aléatoire en revanche,
l'aire de contact suit $p_0$ : les sommets sont chargés dès l'arrivée du pneu, la perturbation dure
plus longtemps et flue davantage. La formule multiplicative sous-estime alors la perturbation
d'environ 9\,\%. En outre, $C=Q/p_0$ est mal conditionné là où $p_0$ est faible. Les résultats du
24/09 sur les textures aléatoires sont donc \textbf{remplacés} par ceux de ce rapport. L'écart entre les deux méthodes est quantifié au \S\ref{sec:verif}.
\end{alerte}

\subsection{Post-traitement}
En chaque point de la grille fine et à chaque profondeur, on assemble le tenseur des
déformations $\bm\epsilon_0+\delta\bm\epsilon$ ; la déformation principale majeure $\epsilon_1$
est sa plus grande valeur propre. On retient deux statistiques sur la zone de la roue arrière
droite :
\begin{itemize}[nosep]
\item le \textbf{maximum}, directement comparable au TFE, mais porté par un seul point de grille ;
\item le \textbf{centile 99,9\,\%}, plus robuste : il est dépassé sur \SI{0.1}{\%} de la surface,
      soit environ \SI{4}{cm^2} sous la roue.
\end{itemize}
Les profondeurs retenues sont 1, 2, 5, 10, 20, 40 et \SI{80}{mm}. Exactement en surface
($z=0$), le modèle de surface plane n'a plus de sens sous une texture, où les pressions sont
quasi ponctuelles : la comparaison avec le TFE se fait donc à partir de \SI{1}{mm}.

%=============================================================================================
\section{Vérifications et degré de confiance}\label{sec:verif}
%=============================================================================================
Chaque maillon de la chaîne a été contrôlé par un calcul indépendant. Le
tableau~\ref{tab:verifs} les résume ; le détail suit.

\begin{table}[H]\centering\small
\caption{Synthèse des vérifications.}\label{tab:verifs}
\begin{tabular}{p{0.8cm}p{4.6cm}p{5.2cm}p{3.6cm}}
\toprule
Test & Ce qui est vérifié & Référence indépendante & Écart obtenu\\
\midrule
V1--V5 & noyau multicouche, charge roulante VE, reproduction du TFE & Love, PyMastic, hérédité, COMSOL & $10^{-4}$ ; < 0,8\,\% ; 0,5--1,5\,\% (voir notice)\\
T1 & contact BEM sur rainures & poinçons périodiques (forme fermée, Legendre) & $2\cdot10^{-5}$ ; coefficients à $10^{-5}$\\
T2 & contact BEM sur texture aléatoire & théorie de Persson & écart connu de 20--30\,\% à faible charge, convergence en contact total\\
T3 & couplage VE complet (hérédité + massif local) & calcul multicouche VE exact, régime \texttt{MovingEnvelope} & 2,3\,\% (empreinte hétérogène, composantes dominantes)\\
T5/T5b & hérédité avec contact non linéaire & marche en temps sur les réponses, deux codes & identiques à 0,00\,\%\\
T6 & massif local contre structure complète & multicouche PEP complet (5 couches) & 0,07\,\% (écart quadratique maximal)\\
T8 & résolution (contact + réponse) & même surface à 0,5 / 0,25 / 0,125 mm & < 1\,\% entre 0,25 et 0,125 mm dès $z\ge\SI{0.5}{mm}$\\
T9 & variabilité d'une réalisation à l'autre & 5 tirages indépendants par MPD & voir \S\ref{sec:tirages}\\
\bottomrule
\end{tabular}
\end{table}

\subsection{Contact : T1 et T2}
Sur les rainures (T1), le BEM converge vers la forme fermée \eqref{eq:punch} (écart de
$6\cdot10^{-4}$, $10^{-4}$ puis $2\cdot10^{-5}$ pour 512, 2048 et 8192 points par pas), et ses
coefficients de Fourier coïncident avec la formule de Legendre à $10^{-5}$ près. Sur les textures
aléatoires (T2), les trois MPD se regroupent sur une même courbe en variable réduite
$\bar p/(E^*\sqrt{m_2})$, comme l'impose la théorie. Persson sous-estime l'aire de 20 à 30\,\% aux
faibles charges, ce qui est l'écart connu de la littérature. \textbf{Le calcul du contact est
fiable} dans son cadre : élasticité linéaire, frottement nul, surface rigide.

\subsection{Couplage deux échelles : T3, T5, T6}
\textbf{T6, massif local.} La réponse à $\delta p$ calculée sur un massif homogène coïncide avec
celle de la structure PEP complète (5 couches, substratum) à mieux que
0,07\,\% en écart quadratique, de 1 à \SI{40}{mm}
de profondeur. L'hypothèse de localisation est donc parfaitement justifiée.

\textbf{T3, viscoélasticité avec texture fixe.} Pour l'harmonique fondamentale du rainurage sous
l'enveloppe hétérogène roulante, la méthode héréditaire \eqref{eq:hered} reproduit le calcul
multicouche viscoélastique exact. Ce dernier est obtenu par un régime spécial de \cs{} où la
pulsation vaut $-\kappa V$ et le noyau est évalué en $\kappa+k_n$. L'écart est au plus de
2,3\,\% sur les composantes dominantes ($\epsilon_{zz}$, $\epsilon_{xz}$, $\epsilon_{xx}$),
et de quelques pour cent sur les composantes faibles. Pour les empreintes en créneau, l'accord reste
de 1 à 4\,\% sur les composantes dominantes.

\textbf{T5, contact non linéaire.} Pour une texture aléatoire, la forme de $\delta p$ change
pendant le passage. On compare deux mises en \oe uvre indépendantes de la même intégrale
d'hérédité. L'une marche en temps sur les \emph{réponses} (t05, une réponse élastique par pas de
temps) ; l'autre, celle de la campagne, applique l'hérédité sur la \emph{pression} (t05b).
\begin{table}[H]\centering\small
\caption{T5b : même intégrale d'hérédité par deux codes (maximum sur une fenêtre de $\SI{64}{mm}\times\SI{64}{mm}$, µdef, perturbation seule).}
\begin{tabular}{llcrr}
\toprule
Texture & Fenêtre & Grandeur & sur les réponses & sur la pression\\
\midrule
gauss MPD1.0 & A & $e_{zz}$ à 2 mm & 390,32 & 390,32 ($-$0,00\,\%)\\
gauss MPD1.0 & A & $e_{xz}$ à 2 mm & 175,22 & 175,22 ($-$0,00\,\%)\\
gauss MPD1.0 & A & $e_{zz}$ à 5 mm & 227,50 & 227,50 ($-$0,00\,\%)\\
gauss MPD1.0 & A & $e_{xz}$ à 5 mm & 105,52 & 105,52 ($-$0,00\,\%)\\
gauss MPD1.0 & B & $e_{zz}$ à 2 mm & 409,46 & 409,46 ($-$0,00\,\%)\\
gauss MPD1.0 & B & $e_{xz}$ à 2 mm & 225,36 & 225,36 ($-$0,00\,\%)\\
gauss MPD1.0 & B & $e_{zz}$ à 5 mm & 236,85 & 236,85 ($-$0,00\,\%)\\
gauss MPD1.0 & B & $e_{xz}$ à 5 mm & 126,85 & 126,85 ($-$0,00\,\%)\\
rainures vives & A & $e_{zz}$ à 2 mm & 167,46 & 167,46 ($-$0,00\,\%)\\
rainures vives & A & $e_{xz}$ à 2 mm & 131,28 & 131,28 ($-$0,00\,\%)\\
rainures vives & A & $e_{zz}$ à 5 mm & 100,95 & 100,95 ($-$0,00\,\%)\\
rainures vives & A & $e_{xz}$ à 5 mm & 61,60 & 61,60 ($-$0,00\,\%)\\
rainures vives & B & $e_{zz}$ à 2 mm & 267,16 & 267,16 ($-$0,00\,\%)\\
rainures vives & B & $e_{xz}$ à 2 mm & 210,32 & 210,32 ($-$0,00\,\%)\\
rainures vives & B & $e_{zz}$ à 5 mm & 161,91 & 161,91 ($-$0,00\,\%)\\
rainures vives & B & $e_{xz}$ à 5 mm & 99,57 & 99,57 ($-$0,00\,\%)\\
\bottomrule
\end{tabular}
\end{table}
Le même test chiffre l'erreur de la formule multiplicative abandonnée (\S\ref{sec:hered}) :
\begin{table}[H]\centering\small
\caption{T5 : méthode héréditaire exacte et formule multiplicative abandonnée (perturbation seule, µdef).}\label{tab:t5}
\begin{tabular}{llcrr}
\toprule
Texture & Fenêtre & Grandeur & héréditaire (exact) & multiplicative\\
\midrule
gauss MPD1.0 & A & $e_{zz}$ à 2 mm & 390,3 & 356,6 ($-$8,6\,\%)\\
gauss MPD1.0 & A & $e_{xz}$ à 2 mm & 175,2 & 160,1 ($-$8,6\,\%)\\
gauss MPD1.0 & A & $e_{zz}$ à 5 mm & 227,5 & 208,8 ($-$8,2\,\%)\\
gauss MPD1.0 & B & $e_{zz}$ à 2 mm & 409,5 & 373,0 ($-$8,9\,\%)\\
gauss MPD1.0 & B & $e_{xz}$ à 2 mm & 225,4 & 208,9 ($-$7,3\,\%)\\
gauss MPD1.0 & B & $e_{zz}$ à 5 mm & 236,9 & 216,4 ($-$8,7\,\%)\\
rainures vives & A & $e_{zz}$ à 2 mm & 167,5 & 167,7 (0,2\,\%)\\
rainures vives & A & $e_{xz}$ à 2 mm & 131,3 & 131,6 (0,2\,\%)\\
rainures vives & A & $e_{zz}$ à 5 mm & 100,9 & 101,1 (0,2\,\%)\\
rainures vives & B & $e_{zz}$ à 2 mm & 267,2 & 267,6 (0,2\,\%)\\
rainures vives & B & $e_{xz}$ à 2 mm & 210,3 & 210,8 (0,2\,\%)\\
rainures vives & B & $e_{zz}$ à 5 mm & 161,9 & 162,2 (0,2\,\%)\\
\bottomrule
\end{tabular}
\end{table}
La formule multiplicative est exacte pour les rainures, et elle \textbf{sous-estime} de 9\,\% environ la perturbation des textures aléatoires. En
effet, sur une texture aléatoire, les sommets sont chargés dès l'arrivée du pneu, avant que
l'enveloppe n'atteigne son maximum : la perturbation dure plus longtemps que ne le suppose la
formule, et le fluage l'amplifie davantage.


\subsection{Convergence en résolution : T8}\label{sec:resolution}
La même surface physique est discrétisée à 0,5, 0,25 et \SI{0.125}{mm}, par troncature spectrale
exacte. Le contact et la réponse sont recalculés sous une pression uniforme de \SI{1.65}{MPa}
(massif élastique de \SI{11670}{MPa}).
\begin{table}[H]\centering\small
\caption{T8 : $\epsilon_1$ maximal (µdef) en fonction du pas de discrétisation.}
\begin{tabular}{lcccccc}
\toprule
Surface & pas (mm) & aire (\%) / $p_{\max}$ & $z$ = 0,5 mm & 1 mm & 2 mm & 5 mm\\
\midrule
aléatoire MPD 1,0 & 0,125 & 54,4 & 177 & 148 & 118 & 78\\
aléatoire MPD 1,0 & 0,250 & 55,4 & 177 & 148 & 118 & 78\\
aléatoire MPD 1,0 & 0,500 & 57,6 & 170 & 147 & 117 & 78\\
rainures, arêtes vives & 0,251 & 7,2 MPa & 121 & 88 & 55 & 15\\
rainures, arêtes vives & 0,125 & 10,2 MPa & 122 & 88 & 56 & 16\\
rainures, arêtes vives & 0,063 & 14,2 MPa & 122 & 88 & 55 & 15\\
rainures, r = 1 mm & 0,251 & 4,3 MPa & 101 & 87 & 60 & 20\\
rainures, r = 1 mm & 0,125 & 4,4 MPa & 103 & 88 & 61 & 20\\
rainures, r = 1 mm & 0,063 & 4,4 MPa & 104 & 88 & 61 & 20\\
\bottomrule
\end{tabular}
\end{table}
Au pas de la campagne (\SI{0.25}{mm}), les grandeurs sont convergées à mieux que 1\,\% dès
\SI{0.5}{mm} de profondeur, y compris pour les arêtes vives. La pression d'arête y est pourtant
singulière (7, 10 puis \SI{14}{MPa} quand le pas diminue), mais son intégrale, seule vue en
profondeur, converge.

\subsection{Degré de confiance : ce qu'on peut affirmer}
\begin{encadre}
\textbf{Fiable (à quelques pour cent).} Dans le cadre des hypothèses (H1)--(H7), les nombres de ce
rapport sont des solutions convergées et vérifiées : le contact, la localisation, la
viscoélasticité avec texture fixe, la résolution et l'accord avec le TFE pour le cas lisse. Les
\textbf{tendances} sont robustes : amplification près de la surface, profondeur d'influence de
10 à \SI{20}{mm}, effet nul en base de couche, classement des textures.

\textbf{Incertain (facteur de l'ordre de 1,5).} La valeur absolue des déformations aux premiers
millimètres dépend fortement de paramètres \emph{mal connus} : le module de la bande de roulement
(voir la sensibilité à $E_r$), la forme exacte de la texture réelle (tirages et asymétrie) et
l'arrondi des arêtes des rainures.

\textbf{Hors de portée du niveau 1, traité par le niveau 2.} La valeur locale dans le mortier
entre granulats (H7), la concentration de contrainte géométrique autour des rainures (H6) et le
frottement (H3). Le niveau 2 (\S\ref{sec:niveau2}) les quantifie : les trois vont dans le sens
d'une sévérité plus grande, comme annoncé. Leur degré de confiance est discuté au
\S\ref{sec:n2_confiance}.
\end{encadre}


%=============================================================================================
\section{Résultats et comparaison avec le TFE}\label{sec:resultats}
%=============================================================================================
Toutes les valeurs concernent la roue arrière droite du bogie, où se trouvent les maxima du TFE,
avec le KVG du TFE à \SI{0.66}{m/s} et \SI{9.3}{\celsius}. Le cas « lisse » est le TFE recalculé
par \cs : il coïncide avec la validation V5 de la notice.

\subsection{Contact}
Les caractéristiques des surfaces et du contact sous \SI{1.65}{MPa} sont données au tableau~\ref{tab:surfaces} (\S\ref{sec:surfaces}).
\begin{figure}[H]\centering\includegraphics[width=0.7\linewidth]{figures/r_aire.pdf}
\caption{Aire de contact en fonction de la pression nominale locale. Sous le pneu, celle-ci varie
de 0 à \SI{3.4}{MPa} (empreinte hétérogène) : la surface en contact change donc fortement d'un
point à l'autre de l'empreinte.}\end{figure}
Sous la pression nominale du TFE, le pneu ne touche que 30 à 85\,\% de la surface selon la
texture et la raideur du caoutchouc. La pression locale y atteint 5 à \SI{15}{MPa}, soit 3 à 9
fois la pression nominale. Sur piste rainurée, le contact couvre les plats (83\,\%) avec une
surpression d'arête.

\subsection{Empreinte hétérogène (cas de référence)}
\begin{table}[H]\centering\footnotesize
\caption{Empreinte hétérogène : déformation principale majeure $\epsilon_1$ maximale (µdef), et entre parenthèses le facteur par rapport au lisse. TFE en surface : 202 µdef (COMSOL), 181 µdef (\cs).}\label{tab:het_e1}
\begin{tabular}{lrrrrrr}
\toprule
Surface \textbackslash{} profondeur & 1 mm & 2 mm & 5 mm & 10 mm & 20 mm & 40 mm\\
\midrule
\textbf{lisse (TFE)} & \textbf{182} & \textbf{184} & \textbf{186} & \textbf{188} & \textbf{185} & \textbf{170}\\ \midrule
rainures FAA, arêtes vives & 317 (1,74) & 305 (1,66) & 234 (1,26) & 197 (1,05) & 185 (1,00) & 169 (1,00)\\
rainures FAA, arêtes $r$ = 1 mm & 315 (1,73) & 306 (1,67) & 243 (1,31) & 203 (1,08) & 186 (1,00) & 169 (1,00)\\
aléatoire, MPD 0,6 mm & 298 (1,64) & 281 (1,53) & 262 (1,41) & 219 (1,17) & 187 (1,01) & 170 (1,00)\\
aléatoire, MPD 1,0 mm & 321 (1,76) & 301 (1,64) & 268 (1,44) & 226 (1,21) & 187 (1,01) & 169 (1,00)\\
aléatoire, MPD 1,5 mm & 358 (1,97) & 327 (1,78) & 271 (1,46) & 231 (1,23) & 187 (1,01) & 169 (1,00)\\
négative, MPD 1,0 mm & 328 (1,80) & 311 (1,70) & 269 (1,45) & 218 (1,16) & 188 (1,01) & 170 (1,00)\\
positive, MPD 1,0 mm & 342 (1,88) & 285 (1,55) & 265 (1,43) & 225 (1,20) & 188 (1,01) & 170 (1,00)\\
MPD 1,0 ; $E_r$ = 5 MPa & 293 (1,61) & 274 (1,49) & 259 (1,39) & 215 (1,15) & 187 (1,01) & 170 (1,00)\\
MPD 1,0 ; $E_r$ = 20 MPa & 395 (2,17) & 342 (1,86) & 272 (1,46) & 236 (1,26) & 188 (1,02) & 170 (1,00)\\
\bottomrule
\end{tabular}
\end{table}

\begin{table}[H]\centering\footnotesize
\caption{Empreinte hétérogène : centile 99,9\,\% de $\epsilon_1$ (µdef), indicateur robuste (dépassé sur environ \SI{4}{cm^2}).}\label{tab:het_e1p}
\begin{tabular}{lrrrrrr}
\toprule
Surface \textbackslash{} profondeur & 1 mm & 2 mm & 5 mm & 10 mm & 20 mm & 40 mm\\
\midrule
\textbf{lisse (TFE)} & \textbf{182} & \textbf{183} & \textbf{185} & \textbf{187} & \textbf{184} & \textbf{169}\\ \midrule
rainures FAA, arêtes vives & 301 (1,66) & 294 (1,60) & 226 (1,22) & 190 (1,02) & 184 (1,00) & 169 (1,00)\\
rainures FAA, arêtes $r$ = 1 mm & 299 (1,64) & 297 (1,62) & 236 (1,27) & 195 (1,04) & 184 (1,00) & 169 (1,00)\\
aléatoire, MPD 0,6 mm & 272 (1,50) & 263 (1,44) & 242 (1,31) & 203 (1,08) & 185 (1,00) & 169 (1,00)\\
aléatoire, MPD 1,0 mm & 291 (1,60) & 281 (1,53) & 253 (1,37) & 211 (1,13) & 185 (1,01) & 169 (1,00)\\
aléatoire, MPD 1,5 mm & 308 (1,69) & 294 (1,60) & 260 (1,40) & 218 (1,17) & 186 (1,01) & 169 (1,00)\\
négative, MPD 1,0 mm & 298 (1,64) & 288 (1,58) & 257 (1,39) & 209 (1,12) & 186 (1,01) & 170 (1,00)\\
positive, MPD 1,0 mm & 279 (1,53) & 262 (1,43) & 247 (1,33) & 209 (1,12) & 186 (1,01) & 169 (1,00)\\
MPD 1,0 ; $E_r$ = 5 MPa & 265 (1,46) & 258 (1,41) & 237 (1,28) & 200 (1,07) & 185 (1,00) & 169 (1,00)\\
MPD 1,0 ; $E_r$ = 20 MPa & 328 (1,80) & 304 (1,66) & 264 (1,43) & 223 (1,19) & 186 (1,01) & 169 (1,00)\\
\bottomrule
\end{tabular}
\end{table}

\begin{table}[H]\centering\footnotesize
\caption{Empreinte hétérogène : $|\epsilon_{xz}|$ maximal (µdef), cisaillement longitudinal. Facteur donné seulement si la valeur lisse dépasse 50 µdef.}\label{tab:het_exz}
\begin{tabular}{lrrrrrr}
\toprule
Surface \textbackslash{} profondeur & 1 mm & 2 mm & 5 mm & 10 mm & 20 mm & 40 mm\\
\midrule
\textbf{lisse (TFE)} & \textbf{8} & \textbf{13} & \textbf{22} & \textbf{32} & \textbf{46} & \textbf{65}\\ \midrule
rainures FAA, arêtes vives & 300 & 221 & 105 & 61 & 51 & 66 (1,01)\\
rainures FAA, arêtes $r$ = 1 mm & 300 & 234 & 127 & 71 & 55 & 66 (1,01)\\
aléatoire, MPD 0,6 mm & 204 & 164 & 100 & 56 & 48 & 66 (1,01)\\
aléatoire, MPD 1,0 mm & 306 & 252 & 149 & 75 & 50 & 66 (1,01)\\
aléatoire, MPD 1,5 mm & 435 & 349 & 200 & 96 & 57 & 66 (1,02)\\
négative, MPD 1,0 mm & 292 & 241 & 145 & 90 & 56 & 64 (0,99)\\
positive, MPD 1,0 mm & 534 & 338 & 164 & 79 & 51 & 66 (1,01)\\
MPD 1,0 ; $E_r$ = 5 MPa & 172 & 140 & 87 & 51 & 48 & 66 (1,01)\\
MPD 1,0 ; $E_r$ = 20 MPa & 574 & 440 & 244 & 113 & 64 & 66 (1,02)\\
\bottomrule
\end{tabular}
\end{table}

\begin{table}[H]\centering\footnotesize
\caption{Empreinte hétérogène : $|\epsilon_{yz}|$ maximal (µdef), c'est-à-dire l'$\epsilon_{xz}$ du TFE (cisaillement transversal). Maximum du TFE en profondeur : 164 µdef vers \SI{100}{mm} (\cs), inchangé par la texture.}\label{tab:het_eyz}
\begin{tabular}{lrrrrrr}
\toprule
Surface \textbackslash{} profondeur & 1 mm & 2 mm & 5 mm & 10 mm & 20 mm & 40 mm\\
\midrule
\textbf{lisse (TFE)} & \textbf{12} & \textbf{23} & \textbf{51} & \textbf{86} & \textbf{125} & \textbf{153}\\ \midrule
rainures FAA, arêtes vives & 14 & 25 & 55 (1,08) & 90 (1,05) & 126 (1,01) & 153 (1,00)\\
rainures FAA, arêtes $r$ = 1 mm & 15 & 26 & 57 (1,10) & 91 (1,06) & 127 (1,02) & 153 (1,00)\\
aléatoire, MPD 0,6 mm & 200 & 179 & 120 (2,33) & 118 (1,37) & 130 (1,04) & 154 (1,00)\\
aléatoire, MPD 1,0 mm & 303 & 275 & 159 (3,10) & 140 (1,64) & 136 (1,09) & 154 (1,01)\\
aléatoire, MPD 1,5 mm & 439 & 376 & 218 (4,25) & 165 (1,92) & 144 (1,15) & 154 (1,01)\\
négative, MPD 1,0 mm & 348 & 277 & 180 (3,51) & 129 (1,51) & 138 (1,10) & 154 (1,01)\\
positive, MPD 1,0 mm & 544 & 351 & 178 (3,46) & 133 (1,55) & 134 (1,08) & 155 (1,01)\\
MPD 1,0 ; $E_r$ = 5 MPa & 172 & 151 & 110 (2,14) & 112 (1,30) & 129 (1,03) & 154 (1,00)\\
MPD 1,0 ; $E_r$ = 20 MPa & 584 & 470 & 276 (5,37) & 187 (2,18) & 150 (1,21) & 154 (1,01)\\
\bottomrule
\end{tabular}
\end{table}

\begin{figure}[H]\centering\includegraphics[width=\linewidth]{figures/r_profils_het.pdf}
\caption{Empreinte hétérogène : maxima des déformations en fonction de la profondeur (échelle
logarithmique). Toutes les courbes rejoignent le cas lisse au-delà de 20 à \SI{40}{mm}.}\label{fig:profils}\end{figure}
\begin{figure}[H]\centering\includegraphics[width=\linewidth]{figures/r_cartes.pdf}
\caption{$\epsilon_1$ à \SI{2}{mm} de profondeur sous la roue arrière droite (empreinte hétérogène).
Lisse : la déformation suit les nervures de l'empreinte. Texturé : elle se fragmente à l'échelle
des sommets de texture et ses pics augmentent.}\end{figure}
\subsection{Variabilité d'une réalisation de texture à l'autre (T9)}\label{sec:tirages}
Une texture aléatoire est une réalisation d'un processus. On répète le calcul complet sur
5 tirages indépendants par MPD : même DSP et même MPD, phases différentes.
\begin{table}[H]\centering\footnotesize
\caption{Facteur d'amplification de $\epsilon_1$ max par rapport au lisse (empreinte hétérogène) : moyenne $\pm$ écart type sur les tirages.}\label{tab:tirages}
\begin{tabular}{lcccccc}
\toprule
Texture & tirages & 1 mm & 2 mm & 5 mm & 10 mm & 20 mm\\
\midrule
MPD 0,6 mm & 5 & 1,62 $\pm$ 0,01 & 1,53 $\pm$ 0,02 & 1,38 $\pm$ 0,04 & 1,15 $\pm$ 0,04 & 1,02 $\pm$ 0,00\\
MPD 1,0 mm & 5 & 1,77 $\pm$ 0,03 & 1,64 $\pm$ 0,03 & 1,43 $\pm$ 0,02 & 1,20 $\pm$ 0,03 & 1,02 $\pm$ 0,00\\
MPD 1,5 mm & 5 & 1,97 $\pm$ 0,02 & 1,75 $\pm$ 0,04 & 1,46 $\pm$ 0,01 & 1,24 $\pm$ 0,03 & 1,02 $\pm$ 0,01\\
\bottomrule
\end{tabular}
\end{table}
\begin{figure}[H]\centering\includegraphics[width=0.66\linewidth]{figures/r_tirages.pdf}
\caption{Amplification moyenne $\pm$ un écart type sur les tirages, en fonction de la profondeur.}\end{figure}
\subsection{Empreintes uniformes}
\begin{table}[H]\centering\footnotesize
\caption{Empreinte rectangulaire : $\epsilon_1$ maximal (µdef) et facteur. TFE en surface : 195 µdef (COMSOL), 196 µdef (\cs).}\label{tab:rect_e1}
\begin{tabular}{lrrrrrr}
\toprule
Surface \textbackslash{} profondeur & 1 mm & 2 mm & 5 mm & 10 mm & 20 mm & 40 mm\\
\midrule
\textbf{lisse (TFE)} & \textbf{215} & \textbf{215} & \textbf{211} & \textbf{203} & \textbf{188} & \textbf{165}\\ \midrule
rainures FAA, arêtes vives & 278 (1,30) & 281 (1,31) & 212 (1,01) & 203 (1,00) & 188 (1,00) & 165 (1,00)\\
rainures FAA, arêtes $r$ = 1 mm & 277 (1,29) & 281 (1,31) & 229 (1,09) & 203 (1,00) & 188 (1,00) & 165 (1,00)\\
aléatoire, MPD 0,6 mm & 300 (1,40) & 284 (1,32) & 248 (1,18) & 204 (1,00) & 188 (1,00) & 165 (1,00)\\
aléatoire, MPD 1,0 mm & 314 (1,46) & 293 (1,36) & 266 (1,26) & 204 (1,01) & 188 (1,00) & 165 (1,00)\\
aléatoire, MPD 1,5 mm & 348 (1,62) & 304 (1,41) & 273 (1,30) & 205 (1,01) & 188 (1,00) & 165 (1,00)\\
négative, MPD 1,0 mm & 310 (1,44) & 295 (1,37) & 273 (1,29) & 221 (1,09) & 189 (1,00) & 165 (1,00)\\
positive, MPD 1,0 mm & 348 (1,62) & 293 (1,36) & 241 (1,15) & 204 (1,01) & 188 (1,00) & 165 (1,00)\\
MPD 1,0 ; $E_r$ = 5 MPa & 292 (1,36) & 279 (1,30) & 237 (1,12) & 203 (1,00) & 188 (1,00) & 165 (1,00)\\
MPD 1,0 ; $E_r$ = 20 MPa & 406 (1,89) & 316 (1,47) & 276 (1,31) & 215 (1,06) & 188 (1,00) & 165 (1,00)\\
\bottomrule
\end{tabular}
\end{table}

\begin{table}[H]\centering\footnotesize
\caption{Empreinte circulaire : $\epsilon_1$ maximal (µdef) et facteur. TFE en surface : 182 µdef (COMSOL), 173 µdef (\cs).}\label{tab:circ_e1}
\begin{tabular}{lrrrrrr}
\toprule
Surface \textbackslash{} profondeur & 1 mm & 2 mm & 5 mm & 10 mm & 20 mm & 40 mm\\
\midrule
\textbf{lisse (TFE)} & \textbf{185} & \textbf{186} & \textbf{183} & \textbf{178} & \textbf{168} & \textbf{150}\\ \midrule
rainures FAA, arêtes vives & 277 (1,50) & 280 (1,50) & 213 (1,16) & 178 (1,00) & 168 (1,00) & 150 (1,00)\\
rainures FAA, arêtes $r$ = 1 mm & 275 (1,49) & 280 (1,50) & 229 (1,25) & 179 (1,00) & 168 (1,00) & 150 (1,00)\\
aléatoire, MPD 0,6 mm & 298 (1,62) & 282 (1,51) & 246 (1,35) & 178 (1,00) & 168 (1,00) & 150 (1,00)\\
aléatoire, MPD 1,0 mm & 312 (1,69) & 291 (1,56) & 263 (1,44) & 187 (1,05) & 168 (1,00) & 149 (1,00)\\
aléatoire, MPD 1,5 mm & 343 (1,86) & 301 (1,62) & 271 (1,48) & 201 (1,13) & 168 (1,00) & 149 (1,00)\\
négative, MPD 1,0 mm & 308 (1,67) & 293 (1,58) & 272 (1,49) & 220 (1,24) & 168 (1,00) & 150 (1,00)\\
positive, MPD 1,0 mm & 343 (1,86) & 291 (1,56) & 240 (1,31) & 178 (1,00) & 168 (1,00) & 149 (1,00)\\
MPD 1,0 ; $E_r$ = 5 MPa & 291 (1,58) & 277 (1,49) & 236 (1,29) & 178 (1,00) & 168 (1,00) & 150 (1,00)\\
MPD 1,0 ; $E_r$ = 20 MPa & 400 (2,16) & 312 (1,68) & 274 (1,50) & 214 (1,20) & 168 (1,00) & 149 (1,00)\\
\bottomrule
\end{tabular}
\end{table}

\subsection{Synthèse et comparaison avec le TFE}
\begin{table}[H]\centering\footnotesize
\caption{Grandeurs clés (µdef). Valeurs du TFE (COMSOL, surface lisse) : $\epsilon_1$ en surface 202 / 195 / 182 et $\epsilon_T$ en base 337 / 331 / 318 µdef (hétérogène / rectangulaire / circulaire). La texture ne modifie pas $\epsilon_T$ en base (effet nul au-delà de \SI{40}{mm}).}\label{tab:synthese}
\begin{tabular}{llrrrrr}
\toprule
Empreinte & Surface & $\epsilon_1$ max & $\epsilon_1$ p99,9 & $\epsilon_1$ max & $|\epsilon_{xz}|$ max & $|\epsilon_{yz}|$ max\\
 & & 2 mm & 2 mm & 10 mm & 2 mm & 40 mm\\
\midrule
hétérogène & lisse (TFE) & 184 & 183 & 188 & 13 & 153\\
 & rainures FAA, arêtes vives & 305 & 294 & 197 & 221 & 153\\
 & aléatoire, MPD 1,0 mm & 301 & 281 & 226 & 252 & 154\\
 & aléatoire, MPD 1,5 mm & 327 & 294 & 231 & 349 & 154\\
\midrule
rectangulaire & lisse (TFE) & 215 & 208 & 203 & 29 & 129\\
 & rainures FAA, arêtes vives & 281 & 278 & 203 & 133 & 129\\
 & aléatoire, MPD 1,0 mm & 293 & 280 & 204 & 252 & 130\\
 & aléatoire, MPD 1,5 mm & 304 & 287 & 205 & 348 & 130\\
\midrule
circulaire & lisse (TFE) & 186 & 177 & 178 & 31 & 110\\
 & rainures FAA, arêtes vives & 280 & 276 & 178 & 137 & 110\\
 & aléatoire, MPD 1,0 mm & 291 & 278 & 187 & 250 & 110\\
 & aléatoire, MPD 1,5 mm & 301 & 284 & 201 & 345 & 111\\
\bottomrule
\end{tabular}
\end{table}
\begin{figure}[H]\centering\includegraphics[width=0.85\linewidth]{figures/r_barres.pdf}
\caption{$\epsilon_1$ maximal à \SI{2}{mm} de profondeur pour les trois empreintes du TFE et les principales surfaces.}\end{figure}
\subsection{Écart avec le bilan du 24 septembre}
\begin{table}[H]\centering\footnotesize
\caption{$\epsilon_1$ max (µdef), empreinte hétérogène : formule multiplicative (bilan du 24/09) / méthode héréditaire (ce rapport).}
\begin{tabular}{lcccc}
\toprule
Surface & 1 mm & 2 mm & 5 mm & 10 mm\\
\midrule
rainures FAA, arêtes vives & 317 / 317 & 305 / 305 & 234 / 234 & 197 / 197\\
aléatoire, MPD 1,0 mm & 333 / 321 & 308 / 301 & 268 / 268 & 226 / 226\\
aléatoire, MPD 1,5 mm & 371 / 358 & 326 / 327 & 270 / 271 & 231 / 231\\
négative, MPD 1,0 mm & 330 / 328 & 311 / 311 & 268 / 269 & 217 / 218\\
\bottomrule
\end{tabular}
\end{table}
Les rainures, où $\delta p$ reste proportionnel à $p_0$, sont quasi inchangées. Pour les
textures aléatoires, les maxima de $\epsilon_1$ ne changent que de quelques pour cent. Deux
effets de la formule multiplicative se compensent en partie : elle sous-estime la perturbation
sous les zones chargées (T5), et elle surestime la souplesse effective dans les zones peu
chargées, où $C=Q/p_0$ est mal conditionné. Les conclusions du bilan restent valables ; les
valeurs de ce rapport les remplacent.

%=============================================================================================
\section{Niveau 2 : frottement, entaille des rainures, granulats et mortier}\label{sec:niveau2}
%=============================================================================================
\newcommand{\NtNeufHomo}{$9,2\times10^{-16}$}
\newcommand{\NtNeufTFE}{$9,1\times10^{-13}$}
\newcommand{\NplanUn}{312}
\newcommand{\NtroisDUn}{317}
\newcommand{\NonzeAgros}{2,7\,\%}
\newcommand{\NonzeAfin}{0,8\,\%}

Le niveau 1 repose sur trois simplifications que le rapport précédent désignait comme ses limites
principales : la chaussée est à surface plane (H6), l'enrobé est homogène à l'échelle du
millimètre (H7) et le contact est sans frottement (H3). Cette section les lève une par une. Le
principe reste le même : chaque modèle ajouté est vérifié contre une solution indépendante, puis
chargé \emph{exactement} par les grandeurs du calcul 3D au point critique, de sorte que son
résultat se raccorde au niveau 1 sans paramètre libre.

\subsection{Frottement pneu--chaussée (levée de H3)}\label{sec:n2_frott}
\paragraph{Modèle.} Au freinage, le pneu exerce sur la chaussée un effort tangentiel dirigé dans
le sens du roulement. On le représente par une loi d'Amontons locale en glissement total :
$q_x(x,y,t)=\mu\,p(x,y,t)$, à l'échelle du pneu \emph{et} à l'échelle de la texture. C'est une
hypothèse majorante pour la texture (en adhérence partielle, le cisaillement local est plus
réparti), et elle est cohérente avec le freinage appuyé, où l'essentiel de l'empreinte glisse. Trois
niveaux : $\mu=0$ (roulement libre, niveau 1), $\mu=0{,}3$ (freinage courant) et $\mu=0{,}6$
(freinage d'urgence sur piste sèche). Par linéarité, la perturbation tangentielle équivalente vaut
$\delta q_{\text{eq}}=\mu\,\delta p_{\text{eq}}$ : la méthode héréditaire (\S\ref{sec:hered}) s'applique
sans modification.

\paragraph{Mise en œuvre.} La réponse nominale avec frottement est calculée par \cs, qui traite
déjà les efforts tangentiels (découplage P-SV / SH). La réponse locale du massif homogène à
$\delta q$ a été ajoutée à \texttt{texture.local\_response}. Pour un massif semi-infini, chaque
amplitude du noyau (normale, tangentielle P-SV, SH) est de la forme $(a+b\,\xi z)\,e^{-\xi z}$
par invariance d'échelle. Les coefficients $(a,b)$ sont identifiés sur le noyau multicouche
général, puis la forme est contrôlée en un troisième point. \textbf{Contrôle T9} : sur un
chargement de texture aléatoire avec $\mu=0{,}5$ et un effort transversal, l'écart à l'évaluation
directe du noyau général vaut \NtNeufHomo{} (massif homogène) et \NtNeufTFE{} (structure complète
du TFE, 5 couches), sur les six composantes et de 1 à \SI{10}{mm}.

\subsection{Modèle local par éléments finis « pixel »}\label{sec:n2_ef}
L'entaille des rainures et la microstructure granulats--mortier sortent du cadre spectral, qui
suppose une couche plane et homogène. Elles sont traitées par un modèle local éléments finis
(\texttt{chausspec/fe2d.py}, écrit pour ce rapport) :
\begin{itemize}[nosep]
\item cellule rectangulaire de largeur $W$, périodique en $x$, de hauteur $H$, encastrée au fond ;
\item maillage régulier de pixels carrés de côté $h$, chacun étant un élément Q4 bilinéaire
      (intégration $2\times2$) portant son propre matériau $(E,\nu)$, $E=0$ désignant un vide ;
\item \emph{déformation plane généralisée} : le déplacement cherché est
      $\bm u=\bm\epsilon^\ast\cdot\bm x+\bm u_c$, où $\bm\epsilon^\ast$ est la déformation nominale
      uniforme ($\epsilon_{xx0}$ dans le plan, $\epsilon_{yy0}$ hors plan) et $\bm u_c$ un champ
      périodique nul au fond. La déformation nominale entre par la précontrainte
      $-\int\bm B^T\bm C:\bm\epsilon^\ast$. Cette formulation « totale » laisse libres les faces
      des vides, ce qu'exige une rainure ;
\item efforts de surface réels (pression et frottement) appliqués sur la face supérieure du
      premier élément solide de chaque colonne ; les fonds de rainure ne sont pas chargés ;
\item $\epsilon_{xy}$ et $\epsilon_{yz}$ nominaux, qui ne sont pas perturbés au premier ordre par
      une géométrie invariante en $y$, sont ajoutés au tenseur avant le calcul de
      $\epsilon_1$ ;
\item résolution directe creuse (factorisation LU unique pour tous les seconds membres).
\end{itemize}

\paragraph{Chargement au point critique.} Pour chaque texture, le calcul 3D (T10) enregistre au
point où $\epsilon_1$ est maximal à \SI{1}{mm} le tenseur nominal complet, la pression
nominale instantanée $p_0$ et la pression nominale héréditaire équivalente
$p_{0,\text{eq}}=E_s\int J(-s)\,\dd_s p_0$. C'est cette dernière qui charge la cellule, modulée par
le contact : $p(x)=p_{0,\text{eq}}\,m(x)$. Dans le cadre (H1)--(H7), cette réduction est exacte
pour la partie texture (\S\ref{sec:hered}). Les déformations nominales $\epsilon_{xx0}$,
$\epsilon_{yy0}$, $\epsilon_{xy0}$ et $\epsilon_{yz0}$ sont prises telles quelles (valeurs
viscoélastiques). \textbf{Conséquence} : la cellule \emph{sans} vide et homogène reproduit le niveau
1. Pour les rainures, qui ne dépendent que de $x$, elle le reproduit exactement : on retrouve
\NplanUn{}~µdef à \SI{1}{mm} contre \NtroisDUn{}~µdef dans le calcul 3D.

\paragraph{Contrôle T11a.} Cellule homogène plane sous la pression de contact des rainures, avec
un frottement $\mu=0{,}5$ et des déformations nominales, comparée à la solution fermée (partie
uniforme analytique, partie périodique par \texttt{local\_response}) : l'écart maximal sur
$\epsilon_{xx}$, $\epsilon_{zz}$ et $\epsilon_{xz}$ vaut \NonzeAgros{} pour $h=\SI{0.2}{mm}$ et
\NonzeAfin{} pour $h=\SI{0.1}{mm}$ à \SI{1}{mm} de profondeur, et moins de
\SI{0.05}{\percent} à \SI{5}{mm}. La convergence est d'ordre 2 environ, comme attendu pour des
éléments Q4. On retient $h=\SI{0.1}{mm}$ (6,35/64).

\paragraph{Distance critique.} Au coin d'une rainure à arêtes vives, la déformation élastique est
singulière : sa valeur au pixel près dépend du maillage et n'a pas de sens physique. On applique
la théorie des distances critiques (Taylor, 2007) : la grandeur pertinente est la déformation
moyennée sur une longueur $L_c$ liée à la microstructure. Pour un enrobé, $L_c$ est de l'ordre de
l'épaisseur des films de mortier et de la taille des plus petits granulats porteurs, soit
\SIrange{0.5}{1}{mm}. Les résultats sont donnés pour $L_c=0$ (pixel), 0,5 et \SI{1}{mm} ; seuls
les deux derniers doivent être interprétés.

\subsection{Entaille des rainures (levée de H6)}\label{sec:n2_rain}
\paragraph{Géométrie.} Un pas de rainurage FAA ($\Lambda=\SI{38.1}{mm}$) avec une rainure
transversale vide de $6{,}35\times\SI{6.35}{mm}$, sur \SI{60}{mm} de profondeur ($W\times H =
384\times605$ pixels de \SI{0.099}{mm}). Deux variantes : arêtes vives, ou arêtes \emph{et} fonds
arrondis de rayon \SI{1}{mm} (arêtes supérieures telles que dans le contact T1, congés de fond
obtenus en remplissant les coins par un quart de cercle). La coupe est dans le plan $(x,z)$ du
roulement, perpendiculaire aux rainures : c'est le plan où la rainure interrompt la couche.
\paragraph{Chargement.} Pression de contact des plateaux (T1, $\bar p=\SI{1.65}{MPa}$) mise à
l'échelle de $p_{0,\text{eq}}$ au point critique « rainures » de T10, frottement éventuel
$\mu p$, déformations nominales du même point. La référence est la même cellule sans vide
(niveau 1). On note $K_{H6}(z)$ le rapport des maxima de $\epsilon_1$ à la profondeur $z$ sous le
plateau.

\subsection{Granulats et mortier (levée de H7)}\label{sec:n2_mortier}
\paragraph{Microstructure.} Coupe 2D de la couche de roulement (\SI{40}{mm} $\times$ \SI{40}{mm},
pixels de \SI{0.1}{mm}) : granulats circulaires de 2 à \SI{10}{mm}, tirés classe par classe du plus
gros au plus petit selon une courbe de Fuller (passant $\propto(d/D)^{0{,}45}$), placés par
tirage séquentiel aléatoire avec un jeu minimal de \SI{0.2}{mm} entre granulats et une
périodicité en $x$. La fraction surfacique obtenue est de 0,53 à 0,56, ce qui est typique des
coupes d'enrobés 0/10 pour la fraction supérieure à \SI{2}{mm}. Le reste est le \emph{mortier
bitumineux} (liant, fines et sables de moins de \SI{2}{mm}), traité comme un milieu homogène. La
surface est sciée : les granulats y sont tronqués. Trois tirages indépendants par cas.
\paragraph{Homogénéisation et mise à l'échelle.} Une cellule biperiodique de même granulométrie,
soumise à trois déformations uniformes, donne le tenseur effectif $\bm C_{\text{hom}}$ dont on
extrait $(E_{\text{hom}},\nu_{\text{hom}})$ isotropes. Les modules des phases sont ensuite mis à
l'échelle pour que $E_{\text{hom}}$ soit le module de l'enrobé équivalent du niveau 1. Seul compte
alors le \emph{contraste} $c=E_{\text{granulat}}/E_{\text{mortier}}$. Avec $\nu_g=0{,}25$ et
$\nu_m=0{,}41$, on obtient $\nu_{\text{hom}}\approx0{,}35$ comme au niveau 1. Le contraste dépend
de la température et du temps de charge : de l'ordre de 10 pour un mortier froid et rapide,
supérieur à 100 pour un mortier chaud sous charge lente (granulats de 50 à \SI{70}{GPa} contre un
mortier de 0,5 à \SI{5}{GPa}). On étudie $c=10$, 30 et 100.
\paragraph{Chargement et indicateurs.} Point critique de la texture aléatoire MPD \SI{1.0}{mm}
(T10) : ligne de pression de contact extraite de la carte 3D à $\bar p=\SI{1.65}{MPa}$ et mise à
l'échelle de $p_{0,\text{eq}}$, plus les déformations nominales. Une ligne de taches de contact
prolongée indéfiniment en $y$ charge plus qu'en 3D : l'amplitude de la modulation est étalonnée
pour que la cellule homogène reproduise $\epsilon_1$ du calcul 3D à \SI{1}{mm} (voir les
résultats ; la ligne brute sert d'étude de sensibilité). On calcule aussi la même cellule
sous une pression \emph{uniforme} $p_{0,\text{eq}}$ (surface lisse, TFE) et sous les seules
déformations nominales. La référence est la cellule homogène $(E_{\text{hom}},\nu_{\text{hom}})$
sous le même chargement. Indicateur : facteur de concentration dans le mortier
$K_m(z)=\max\epsilon_1^{\text{mortier}}/\max\epsilon_1^{\text{homogène}}$ dans une bande de
$\pm\SI{0.25}{mm}$ autour de $z$ (maximum et quantile à 99\,\%).

\subsection{Résultats : frottement}

\begin{figure}[H]\centering
\includegraphics[width=\textwidth]{figures/n2_frottement.pdf}
\caption{Déformation principale majeure maximale en fonction de la profondeur, empreinte
hétérogène, pour trois niveaux de frottement (calcul 3D complet, T10). Échelle de profondeur
logarithmique.}\label{fig:n2_frott}
\end{figure}

Le frottement augmente $\epsilon_1$ à toutes les profondeurs superficielles
(figure~\ref{fig:n2_frott}, tableau~\ref{tab:n2_frott}). Surface lisse, à \SI{1}{mm} :
182, 202 puis 266~µdef pour $\mu$ = 0, 0,3 et 0,6, soit
$\times$1,11 et $\times$1,46. L'effet est \emph{plus fort sur
les surfaces texturées}, parce que le frottement se concentre lui aussi sur les sommets : à
\SI{1}{mm}, les rainures passent de 317 à 507~µdef ($\times$1,60)
et la macrotexture MPD \SI{1.5}{mm} de 358 à 681~µdef
($\times$1,90). Le rapport texturé/lisse, qui résume l'effet de la texture,
\emph{augmente} donc avec le frottement : pour la MPD \SI{1.5}{mm} à \SI{1}{mm}, il passe de
1,97 sans frottement à 2,56 avec $\mu=0{,}6$. En profondeur
(\SI{10}{mm} et au-delà), le frottement agit par l'échelle du pneu et de la même façon sur
toutes les surfaces ; la conclusion du niveau 1 (pas d'effet de la texture au-delà de 20 à
\SI{40}{mm}) est inchangée. Le frottement modifie en revanche les grandeurs du TFE
elles-mêmes, indépendamment de la texture : sur surface lisse, $\epsilon_1$ augmente encore de
29\,\% à \SI{10}{mm} avec $\mu=0{,}6$. Le freinage
(et le virage) mériterait donc une étude propre à l'échelle du TFE.

\begin{table}[H]\centering\small
\caption{$\epsilon_1$ maximale (µdef) avec frottement ; entre parenthèses, le rapport au cas sans
frottement. Deux dernières colonnes : rapport texturé/lisse sans frottement et avec $\mu=0{,}6$.}\label{tab:n2_frott}
\begin{tabular}{llrrrrr}\toprule
surface & $z$ (mm) & $\mu=0$ & $\mu=0{,}3$ & $\mu=0{,}6$ & tex./lisse, $\mu=0$ & tex./lisse, $\mu=0{,}6$\\\midrule
lisse (TFE) & 1 & 182 & 202 (1,11) & 266 (1,46) & 1,00 & 1,00\\
 & 2 & 184 & 205 (1,12) & 252 (1,37) & 1,00 & 1,00\\
 & 5 & 186 & 208 (1,12) & 247 (1,33) & 1,00 & 1,00\\
 & 10 & 188 & 200 (1,07) & 242 (1,29) & 1,00 & 1,00\\
\addlinespace
rainures FAA & 1 & 317 & 407 (1,29) & 507 (1,60) & 1,74 & 1,91\\
 & 2 & 306 & 362 (1,18) & 459 (1,50) & 1,66 & 1,82\\
 & 5 & 234 & 252 (1,08) & 348 (1,48) & 1,26 & 1,41\\
 & 10 & 197 & 219 (1,11) & 260 (1,32) & 1,05 & 1,08\\
\addlinespace
aléatoire MPD 1,0 mm & 1 & 321 & 364 (1,13) & 549 (1,71) & 1,76 & 2,07\\
 & 2 & 301 & 332 (1,10) & 436 (1,45) & 1,64 & 1,73\\
 & 5 & 268 & 287 (1,07) & 336 (1,25) & 1,44 & 1,36\\
 & 10 & 226 & 244 (1,08) & 282 (1,25) & 1,21 & 1,17\\
\addlinespace
aléatoire MPD 1,5 mm & 1 & 358 & 471 (1,32) & 681 (1,90) & 1,97 & 2,56\\
 & 2 & 327 & 363 (1,11) & 528 (1,62) & 1,78 & 2,09\\
 & 5 & 271 & 295 (1,09) & 371 (1,37) & 1,46 & 1,50\\
 & 10 & 231 & 250 (1,08) & 287 (1,24) & 1,23 & 1,19\\
\bottomrule\end{tabular}\end{table}

\subsection{Résultats : entaille des rainures}

\begin{figure}[H]\centering
\includegraphics[width=\textwidth]{figures/n2_rainures_cartes.pdf}
\caption{Déformation principale majeure (moyennée sur $L_c$ = 0,5 mm) dans les 15 premiers
millimètres d'un pas de rainurage, au point critique, sans frottement. La rainure est centrée ;
la pression de contact n'est appliquée que sur les plateaux.}\label{fig:n2_rain_cartes}
\end{figure}

\begin{figure}[H]\centering
\includegraphics[width=\textwidth]{figures/n2_rainures_profils.pdf}
\caption{À gauche : facteur d'entaille $K_{H6}(z)$ (rapport des maxima de $\epsilon_1$ à la
profondeur $z$, avec et sans rainure, $L_c$ = 0,5 mm). À droite : convergence en maillage du
maximum global et du maximum à 1 mm, pour trois longueurs de moyenne.}\label{fig:n2_rain_prof}
\end{figure}

\paragraph{Ce que montre le calcul.} La rainure modifie fortement le champ des premiers
millimètres (figure~\ref{fig:n2_rain_cartes}). Deux zones se distinguent :
\begin{itemize}[nosep]
\item \textbf{L'arête supérieure du plateau.} La surpression de contact au bord (T1) et la face
      libre de la rainure s'y combinent : sous la surface, le matériau n'est plus confiné
      latéralement et s'étend vers la rainure. C'est là que se trouve le maximum global :
      1010~µdef à $z$ = 0,05~mm au maillage de référence (valeur non convergée
      pour une arête parfaitement vive, voir plus bas) ($L_c$ = 0,5 mm ; 898~µdef
      avec $L_c$ = 1 mm), contre 324~µdef au maximum dans la cellule plane.
\item \textbf{Le fond et les coins inférieurs de la rainure}, à \SI{6.35}{mm} : la
      déformation nominale, qui ne peut plus passer par la partie évidée, se concentre autour du
      fond. On y trouve 740~µdef, contre 163 à la même profondeur
      dans la cellule plane ($K_{H6}$ = 4,5). C'est le second point chaud, un peu
      moins sollicité que l'arête pour des arêtes vives.
\end{itemize}
Le facteur d'entaille vaut $K_{H6}$ = 2,11 à \SI{1}{mm},
1,70 à \SI{2}{mm} et 2,44 à \SI{5}{mm}. Il vaut encore
2,15 à \SI{10}{mm}, soit \SI{3.65}{mm} sous le fond de rainure, et
1,18 à \SI{20}{mm} (figure~\ref{fig:n2_rain_prof}) : la profondeur d'influence
de l'entaille est fixée par la profondeur de la rainure, et non plus par la seule taille des
zones de contact. Au-delà de \SI{5}{mm}, ce facteur ne doit pas être appliqué au maximum du
calcul 3D, qui se trouve ailleurs dans l'empreinte ; il décrit l'état local sous la rainure.

\paragraph{Décomposition.} Sous les seules déformations nominales, la rainure donne
$K$ = 0,62 à \SI{1}{mm} (maximum 507~µdef à 6,40~mm, contre
284~µdef en cellule plane) ; sous la seule pression de contact, $K$ = 2,35
(maximum 898~µdef à 0,05~mm, contre 285). La
déformation nominale seule est même \emph{soulagée} sur le plateau à \SI{1}{mm} (la rainure
libère la compression horizontale), mais elle est concentrée au fond. C'est la pression de
contact, concentrée au bord du plateau libre, qui fait le maximum près de la surface.

\paragraph{Arrondis et frottement.} Des arêtes et des fonds arrondis à \SI{1}{mm} réduisent le
maximum global de 1010 à 825~µdef et le facteur à \SI{1}{mm} de
2,11 à 1,59. Le point chaud se déplace alors au fond de la
rainure ($z$ = 6,2~mm), où l'arrondi de \SI{1}{mm} ne suffit pas à faire disparaître la
concentration : 727~µdef à \SI{6.35}{mm}. L'arrondi des arêtes, que l'usure et le
trafic produisent naturellement, est donc bénéfique près de la surface mais ne supprime pas
l'effet. Avec frottement, $K_{H6}$ à \SI{1}{mm}
vaut 2,00 ($\mu=0{,}3$) et 1,92 ($\mu=0{,}6$) pour des arêtes vives ;
avec des arêtes arrondies, le maximum global atteint 1132 ($\mu=0{,}3$) puis
1457~µdef ($\mu=0{,}6$), contre 825~µdef sans frottement.

\paragraph{Robustesse numérique.} L'étude de maillage (cellule de \SI{30}{mm} de profondeur,
pour tenir en mémoire au maillage le plus fin ; $h$ = 0,198 ; 0,099 ; 0,066~mm)
sépare deux types de grandeurs.
\begin{itemize}[nosep]
\item \textbf{Stables à quelques pour cent} : le maximum à \SI{1}{mm} (619,
      658 puis 680~µdef, soit
      $K_{H6}(1\,\text{mm})$ de 2,00 à
      2,18), et plus généralement tout ce qui
      est à plus d'un demi-millimètre des arêtes.
\item \textbf{Non convergé} : le maximum global au coin supérieur d'une arête vive
      (840, 1010 puis
      1131~µdef avec $L_c$ = 0,5 mm). Il est piloté par la surpression de
      contact à l'arête, elle-même singulière pour une arête parfaitement vive (T1). Ce maximum
      n'a pas de valeur physique : une arête réelle est toujours émoussée. On retient la variante
      arrondie, qui supprime la singularité de contact.
\end{itemize}

\paragraph{Limite propre à ce calcul.} La coupe est 2D : elle représente une rainure infiniment
longue sous une déformation hors plan imposée. C'est exact pour la géométrie (les rainures sont
continues sur la largeur de la piste) mais approché pour $\epsilon_{yz}$, qui n'est pas perturbé
dans la cellule ; au point critique, $\epsilon_{yz}$ nominal est faible
(0,1~µdef), et l'erreur induite est donc négligeable.

\subsection{Résultats : granulats et mortier}

\begin{figure}[H]\centering
\includegraphics[width=\textwidth]{figures/n2_mortier_cartes.pdf}
\caption{Déformation principale majeure dans les 15 premiers millimètres (fenêtre de 20 mm de
large, à l'échelle), sous la pression de contact de la texture MPD 1,0 mm au point critique : enrobé homogène équivalent (niveau 1) et
microstructure granulats--mortier (premier tirage) pour deux contrastes de module $c$. Contours
blancs : granulats.}\label{fig:n2_mort_cartes}
\end{figure}

\begin{figure}[H]\centering
\includegraphics[width=\textwidth]{figures/n2_mortier_profils.pdf}
\caption{Facteur de concentration dans le mortier $K_m(z)$ sous la texture (à gauche) et sous une
surface lisse (au centre) : maximum (trait plein, bande = min--max sur 3 tirages) et quantile à
99\,\% (tirets). À droite : rapport $K_m^{\text{texture}}/K_m^{\text{lisse}}$ (maximum en trait
plein, quantile à 99\,\% en tirets) ; une valeur inférieure à 1 signifie que la microstructure
atténue l'écart relatif entre texture et surface lisse.}\label{fig:n2_mort_prof}
\end{figure}

\begin{table}[H]\centering\small
\caption{Granulats et mortier (moyenne de 3 tirages, chargement étalonné). $K_m$ : rapport du
maximum de $\epsilon_1$ dans le mortier au maximum dans l'enrobé homogène équivalent, à la même
profondeur et sous le même chargement. Trois dernières colonnes : $K_m$ texture / $K_m$ lisse.}\label{tab:n2_mort}
\begin{tabular}{rrr ccc ccc ccc}\toprule
& & & \multicolumn{3}{c}{$K_m$, texture} & \multicolumn{3}{c}{$K_m$, lisse} & \multicolumn{3}{c}{rapport}\\
\cmidrule(lr){4-6}\cmidrule(lr){7-9}\cmidrule(lr){10-12}
$c$ & fraction & $E_{\text{hom}}/E_m$ & 1 mm & 2 mm & 5 mm & 1 mm & 2 mm & 5 mm & 1 mm & 2 mm & 5 mm\\\midrule
10 & 0,54 & 2,84 & 2,74 & 2,95 & 2,77 & 3,55 & 4,27 & 4,31 & 0,77 & 0,71 & 0,64\\
30 & 0,54 & 3,56 & 3,29 & 3,86 & 3,29 & 4,85 & 5,60 & 5,51 & 0,70 & 0,71 & 0,60\\
100 & 0,54 & 3,95 & 3,65 & 4,38 & 3,64 & 5,83 & 6,36 & 6,36 & 0,65 & 0,71 & 0,58\\
\bottomrule\end{tabular}\end{table}

\paragraph{Ce que montre le calcul.} Dans un enrobé réel, le mortier se déforme beaucoup plus
que la moyenne : les granulats, 10 à 100 fois plus raides, ne se déforment presque pas, et la
déformation se reporte sur les films de mortier qui les séparent
(figure~\ref{fig:n2_mort_cartes}). Sous la texture, le facteur de concentration $K_m$ vaut
2,7 à 3,7 à \SI{1}{mm} selon le contraste
(tableau~\ref{tab:n2_mort}) ; le quantile à 99\,\%, moins sensible aux films les plus minces, vaut
2,3 à 3,1. $K_m$ croît avec le contraste,
mais de moins en moins : au-delà de $c\approx30$, les granulats sont pratiquement indéformables.

\paragraph{Cette concentration existe aussi sous une surface lisse, et elle y est plus forte.} Sous
une pression uniforme, $K_m$ vaut 3,6 à 5,8 à \SI{1}{mm}
et reste du même ordre à \SI{10}{mm} (3,3 à 5,5).
C'est une propriété du matériau, et non de la texture : elle concerne aussi le TFE, à toutes les
profondeurs. Les lois de fatigue, calées sur des déformations moyennes d'enrobé, la contiennent
implicitement ; il ne faut pas l'ajouter une seconde fois à une vérification en déformation moyenne.

\paragraph{Conséquence : dans le mortier, l'effet \emph{relatif} de la texture est atténué.} Le rapport
$K_m^{\text{texture}}/K_m^{\text{lisse}}$ vaut 0,65 à 0,77 à \SI{1}{mm}
(figure~\ref{fig:n2_mort_prof}, à droite). Sous une surface lisse, la déformation nominale se
concentre dans tous les films de mortier ; sous la texture, le maximum est imposé par les
sommets chargés, qui s'appuient souvent sur des granulats affleurants. Dans le mortier, le rapport
texturé/lisse du niveau 1 (1,76 à \SI{1}{mm} pour la MPD \SI{1.0}{mm}) devient
donc environ 1,1 à 1,4. En valeur absolue, en
revanche, la déformation dans le mortier sous une texture est bien plus forte que la déformation
moyenne du niveau 1 : $\epsilon_1\approx$ 878 à 1173~µdef à
\SI{1}{mm} au lieu de 321. Les écarts entre tirages (bandes de la
figure~\ref{fig:n2_mort_prof}) sont importants sur le maximum ponctuel, qui dépend de la position d'un
film mince sous un sommet de texture, et plus faibles sur le quantile à 99\,\%.

\paragraph{Sensibilité au chargement 2D.} Une ligne de pression extraite de la carte de contact 3D
et prolongée indéfiniment dans la coupe (déformation plane) charge davantage que les taches de
contact réelles. La modulation a donc été étalonnée, $m_{\text{cal}}=1+\alpha(m-1)$, pour que la
cellule homogène reproduise $\epsilon_1$ du calcul 3D au point critique à \SI{1}{mm}
(321~µdef ; la ligne brute donne 375~µdef) :
$\alpha$ = 0,82. Avec la ligne brute, $K_m$ à \SI{1}{mm} vaut
2,7 à 3,7 : le facteur de concentration est peu sensible à la sévérité du
chargement, ce qui rend le résultat robuste à cet étalonnage.

\paragraph{Limites propres à ce calcul.} Coupe 2D (les films de mortier y sont plus continus
qu'en 3D, ce qui majore un peu $K_m$), granulats circulaires (les granulats concassés anguleux
concentrent davantage aux arêtes), mortier élastique au module du point de fonctionnement, et
interface granulat--mortier parfaite (pas de décohésion).
\subsection{Synthèse : résultats corrigés}\label{sec:n2_synthese}

Le tableau~\ref{tab:n2_synthese} rassemble les déformations de l'enrobé équivalent
(déformation moyenne, celle que vérifient les méthodes de dimensionnement), corrigées du
frottement (calcul 3D, T10) et, pour les rainures, de l'entaille (T11 : fourchette entre $L_c$ =
0,5 et 1 mm et, sans frottement, entre arêtes vives et arrondies). La déformation locale dans
le mortier s'en déduit en multipliant par $K_m$ (tableau~\ref{tab:n2_mort}).
\begin{table}[H]\centering\small
\caption{$\epsilon_1$ maximale (µdef) dans l'enrobé équivalent, empreinte hétérogène du TFE.
Niveau 1 : sans frottement, surface plane. Niveau 2 : avec frottement $\mu$ et, pour les rainures,
entaille géométrique. Trois dernières colonnes : rapport à la surface lisse dans les mêmes
conditions.}\label{tab:n2_synthese}
\begin{tabular}{llr ccc ccc}\toprule
& & niveau 1 & \multicolumn{3}{c}{niveau 2} & \multicolumn{3}{c}{rapport à la surface lisse}\\
\cmidrule(lr){4-6}\cmidrule(lr){7-9}
surface & $z$ (mm) & $\mu=0$ & $\mu=0$ & $\mu=0{,}3$ & $\mu=0{,}6$ & niveau 1 & niv. 2, $\mu=0$ & niv. 2, $\mu=0{,}6$\\\midrule
lisse (TFE) & 1 & 182 & 182 & 202 & 266 & 1,00 & 1,00 & 1,00\\
 & 2 & 184 & 184 & 205 & 252 & 1,00 & 1,00 & 1,00\\
 & 5 & 186 & 186 & 208 & 247 & 1,00 & 1,00 & 1,00\\
\addlinespace
rainures FAA & 1 & 317 & 495--668 & 750--816 & 904--975 & 1,74 & 2,7--3,7 & 3,4--3,7\\
 & 2 & 306 & 406--519 & 624--645 & 752--785 & 1,66 & 2,2--2,8 & 3,0--3,1\\
 & 5 & 234 & 502--572 & 615--625 & 734--744 & 1,26 & 2,7--3,1 & 3,0--3,0\\
\addlinespace
aléatoire MPD 1,0 mm & 1 & 321 & 321 & 364 & 549 & 1,76 & 1,76 & 2,07\\
 & 2 & 301 & 301 & 332 & 436 & 1,64 & 1,64 & 1,73\\
 & 5 & 268 & 268 & 287 & 336 & 1,44 & 1,44 & 1,36\\
\addlinespace
aléatoire MPD 1,5 mm & 1 & 358 & 358 & 471 & 681 & 1,97 & 1,97 & 2,56\\
 & 2 & 327 & 327 & 363 & 528 & 1,78 & 1,78 & 2,09\\
 & 5 & 271 & 271 & 295 & 371 & 1,46 & 1,46 & 1,50\\
\bottomrule\end{tabular}\end{table}
Lecture :
\begin{itemize}[nosep]
\item \textbf{Surface lisse (TFE).} Le frottement seul ajoute 11 à
      46\,\% à \SI{1}{mm}.
\item \textbf{Rainures.} C'est la surface la plus modifiée par le niveau 2 : l'entaille et le
      frottement se cumulent. Avec $\mu=0{,}3$, $\epsilon_1$ à \SI{1}{mm} vaut 750 à 816~µdef,
      soit 3,7 à 4,0 fois la
      surface lisse, contre 1,74 au niveau 1.
      Le niveau 1 sous-estimait donc nettement la sévérité des rainures, comme il l'annonçait
      (« borne inférieure »).
\item \textbf{Macrotextures aléatoires.} Seul le frottement les modifie dans ce tableau (pas
      d'entaille) : le rapport à la surface lisse passe de 1,97 à
      2,56 pour la MPD \SI{1.5}{mm} à \SI{1}{mm}. Dans le
      mortier, les déformations absolues sont 2,7 à 3,7 fois
      plus fortes, mais l'écart relatif entre texture et surface lisse est atténué
      (tableau~\ref{tab:n2_mort}).
\item \textbf{En profondeur}, la conclusion du niveau 1 tient : à \SI{40}{mm}, les écarts entre
      surfaces restent nuls (moins de 0,3\,\%) avec ou sans frottement. À \SI{20}{mm}, ils
      atteignent 5 à 7\,\% avec $\mu=0{,}6$, contre 1 à 2\,\% sans frottement
      (figure~\ref{fig:n2_frott}).
\end{itemize}

\subsection{Degré de confiance du niveau 2}\label{sec:n2_confiance}

\begin{encadre}
\textbf{Fiable.} Le solveur EF est vérifié contre la solution fermée (T11a, moins de 1\,\% à
\SI{1}{mm}), la cellule plane reproduit le calcul 3D, et la réponse au frottement est vérifiée à
la précision machine (T9). Les \emph{sens} de variation sont sûrs : frottement, entaille et
microstructure augmentent tous les trois la déformation superficielle, et l'effet de la texture
reste confiné aux 10 à \SI{20}{mm} supérieurs. Pour la microstructure, le passage de 0,1 à \SI{0.067}{mm} (contraste 30, premier tirage)
change $K_m$ à \SI{1}{mm} de 3,06 à 3,13 (maximum) et de
2,73 à 2,76 (quantile à 99\,\%), et le rapport
$K_m^{\text{texture}}/K_m^{\text{lisse}}$ de 0,58 à 0,60.

\textbf{Semi-quantitatif (facteur 1,5 à 2).} Les facteurs du niveau 2 dépendent de paramètres
que seules des mesures fixeront : la distance critique $L_c$ et l'arrondi réel des arêtes pour
les rainures ; le contraste granulat/mortier et la forme des granulats pour $K_m$ ; le coefficient
de frottement mobilisé. Les fourchettes des tableaux couvrent ces incertitudes.

\textbf{À ne pas faire.} Multiplier les résultats du TFE par $K_m$ pour une vérification en
fatigue classique : les lois de fatigue sont calées sur des déformations moyennes d'enrobé, qui
contiennent déjà implicitement la concentration dans le mortier. $K_m$ sert à comparer des
situations entre elles (texturé contre lisse, rainuré contre plan) et, à terme, à un critère
d'amorçage calé sur essais de mortier.
\end{encadre}


%=============================================================================================
\section{Discussion, limites et suite}\label{sec:discussion}
%=============================================================================================
\subsection{Interprétation mécanique}
Le pneu ne porte que sur une fraction de la surface : 30 à 85\,\% selon la texture. La pression
réelle est donc concentrée sur les sommets, à 3 à 9 fois la pression nominale. Cette surpression
se diffuse dans le massif sur une profondeur de l'ordre de la taille des zones de contact, soit
quelques millimètres, et la longueur d'onde de la texture fixe la profondeur d'influence.
Au-delà, seule compte la résultante, identique au cas lisse (principe de Saint-Venant). C'est
pourquoi le TFE, qui a montré une sensibilité croissante à l'empreinte quand on se rapproche de
la surface, trouve ici un prolongement naturel : à l'échelle du pneu, la forme de l'empreinte
agit sur les premiers centimètres ; à l'échelle de la texture, les sommets agissent sur les
premiers millimètres.

La viscoélasticité joue un rôle non trivial. Les sommets de texture sont chargés dès l'arrivée
du pneu et pendant tout son passage, et la roue avant laisse une déformation différée au même
endroit de la chaussée avant le passage de la roue arrière. La méthode héréditaire en tient
compte exactement ; l'ancienne formule multiplicative sous-estimait l'effet des textures aléatoires d'environ 9\,\%.

\subsection{Limites}
Les trois limites structurelles du niveau 1 (H3, H6, H7) sont levées par le niveau 2 ; celles qui
restent sont les suivantes.
\begin{enumerate}[nosep]
\item \textbf{Caoutchouc.} Le module de la bande de roulement du pneu d'A340 (et sa
      dépendance à la fréquence et à la température) reste le paramètre le plus influent et le
      moins connu. Données à obtenir auprès du manufacturier ou dans la littérature (Michelin,
      Airbus).
\item \textbf{Textures synthétiques.} Elles sont réalistes dans leurs statistiques mais ne
      remplacent pas des relevés réels. Le code accepte directement un relevé de hauteurs (laser
      3D) à la place de la surface générée.
\item \textbf{Modèles locaux 2D.} L'entaille et la microstructure sont traitées en coupe
      (déformation plane généralisée), avec des granulats circulaires et un mortier élastique au
      point de fonctionnement. Une cellule 3D (granulats anguleux, mortier viscoélastique) est
      l'étape suivante naturelle ; le solveur EF s'étend au 3D sans changement de principe.
\item \textbf{Frottement} en glissement total (loi d'Amontons locale) : majorant pour la
      texture, et représentatif d'un freinage appuyé. Le virage (effort transversal) n'a pas été
      calculé ; \cs{} le traite déjà.
\item \textbf{Un seul cas de chargement} : vitesse et température du TFE. À vitesse plus élevée
      ou à température plus basse, l'enrobé est plus raide : la déformation diminue, mais pas
      nécessairement la contrainte, et le contraste granulat/mortier diminue.
\item \textbf{Critère d'endommagement.} Le rapport compare des déformations. Relier la
      déformation locale dans le mortier à l'amorçage (fissuration par le haut, arrachement)
      demande un critère calé sur essais (mortier bitumineux, essais de fatigue de mortier).
\end{enumerate}

\subsection{Suite proposée}
\begin{enumerate}[nosep]
\item \textbf{Relevés de texture réels} sur piste (STAC, profilomètre laser), calage de la DSP et
      comparaison avec les surfaces synthétiques.
\item \textbf{Module de la bande de roulement} : données pneumatiques, ou identification inverse
      par des mesures de pression de contact sur surface texturée (capteurs en film).
\item \textbf{Cellule mésoscopique 3D} à partir d'une tomographie d'enrobé réel, mortier
      viscoélastique ; comparaison avec la cellule 2D pour chiffrer l'effet de coupe.
\item \textbf{Essais sur mortier bitumineux} (module complexe et fatigue) pour transformer
      $K_m$ en critère d'amorçage superficiel.
\item \textbf{Indicateur mécanique de texture} : relier l'amplification à MPD, $m_2$ et $S_{sk}$
      sur une famille de surfaces, et proposer un indicateur utilisable en ingénierie.
\item \textbf{Confrontation aux dégradations observées} : fissuration par le haut et arrachement
      aux bords de rainures (Toulouse-Blagnac).
\end{enumerate}


\appendix
\section{Paramètres numériques et temps de calcul}\label{ann:param}
\begin{center}\small
\begin{tabular}{p{6.2cm}p{8.3cm}}
\toprule
Cellule de texture & $\SI{128}{mm}\times\SI{128}{mm}$, $512^2$ points, pas \SI{0.25}{mm}, périodique\\
DSP & $H=0{,}8$ ; $\lambda\in[1,50]$ mm ; plateau au-delà de \SI{20}{mm}\\
Bande de roulement & $E_r=10$ (5 ; 20) MPa, $\nu_r=0{,}49$, $t_r=\SI{10}{mm}$ sur ceinture rigide\\
Contact & gradient conjugué projeté, tolérance $10^{-7}$, 8 niveaux de $\bar p$ de 0,15 à \SI{4}{MPa}\\
Réponse nominale & \cs, domaine $6\times3$ m, $2048\times1024$ (2,9 mm), KVG, $V=\SI{0.66}{m/s}$\\
Fenêtre fine & roue arrière droite, $x\in[-2{,}30;-1{,}66]$ m, $y\in[0{,}40;1{,}00]$ m, marge \SI{2}{cm}, pas \SI{0.25}{mm}\\
Histoire & depuis l'arrivée de la roue avant, pas de 10 mm (hétérogène) ou 5 mm (créneaux)\\
Temps par cas & bibliothèque de contact 40 s ; nominal 100 s ; texture 3 à 6 min\\
\addlinespace
\multicolumn{2}{l}{\emph{Niveau 2}}\\
Frottement (T10) & $\mu$ = 0 ; 0,3 ; 0,6, glissement total, pneu et texture ; 4 surfaces\\
Cellule rainure (T11) & $38{,}1\times\SI{60}{mm}$, pixels de \SI{0.099}{mm} ($384\times605$), Q4, fond encastré ; 20 à 40 s par cas\\
Cellule granulats (T12) & $40\times\SI{40}{mm}$, pixels de \SI{0.1}{mm}, granulats 2--\SI{10}{mm} (Fuller 0,45), jeu \SI{0.2}{mm}, 3 tirages, $c$ = 10 ; 30 ; 100\\
Distance critique & $L_c$ = 0 ; 0,5 ; \SI{1}{mm} (moyenne glissante sur les seuls pixels solides)\\
\bottomrule
\end{tabular}
\end{center}

\section{Reproduction}\label{ann:repro}
Depuis le dossier \texttt{chausspec} (Python 3, NumPy, SciPy, Matplotlib) :
\begin{verbatim}
python validation/t01_t02_contact.py          # T1-T2 : contact
python validation/t03_deux_echelles.py        # T3 : couplage / multicouche exact
python validation/t05_integration_temporelle.py
python validation/t05b_controle_hereditaire.py  # T5 : hérédité, deux codes
python validation/t06_demi_espace_local.py    # T6 : massif local
python validation/t08_resolution.py           # T8 : résolution
ONLY=heterogene OUT=_het \
    python validation/t07_campagne_definitive.py
ONLY=rectangulaire,circulaire OUT=_rc \
    python validation/t07_campagne_definitive.py
SEEDS=2,3,4,5 ONLY=heterogene OUT=_seeds \
    python validation/t07_campagne_definitive.py
python validation/t09_tangentiel_local.py      # T9 : frottement, réponse locale
for mu in 0 0.3 0.6; do ONLY=heterogene \
  TEX=rainures_vives,gauss_MPD1.0,gauss_MPD1.5 MU=$mu OUT=_mu$mu \
  python validation/t10_frottement.py; done     # T10 : frottement 3D
python validation/t11a_ef_validation.py       # T11a : solveur EF
python validation/t11_entaille_rainures.py    # T11 : entaille des rainures
python validation/t12_granulats_mortier.py    # T12 : granulats / mortier
python rapport/figures_methodo.py
python rapport/figures_resultats.py
python rapport/figures_niveau2.py
python rapport/texte_niveau2.py
pdflatex rapport/rapport_texture.tex
\end{verbatim}

\section*{Références}
\small
\begin{itemize}[nosep,leftmargin=*]
\item FAA (2016). AC 150/5320-12C, \emph{Measurement, construction and maintenance of skid-resistant airport pavement surfaces}.
\item OACI. Annexe 14, Vol.~I --- texture de surface des pistes.
\item ISO 13473-1. Caractérisation de la texture d'un revêtement à partir de relevés de profils --- Partie 1 : profondeur moyenne de profil.
\item Johnson K.L. (1985). \emph{Contact Mechanics}. Cambridge University Press.
\item Persson B.N.J. (2001). Theory of rubber friction and contact mechanics. \emph{J. Chem. Phys.} 115, 3840--3861.
\item Polonsky I.A., Keer L.M. (1999). A numerical method for solving rough contact problems based on the multi-level multi-summation and conjugate gradient techniques. \emph{Wear} 231, 206--219.
\item Taylor D. (2007). \emph{The Theory of Critical Distances: A New Perspective in Fracture Mechanics}. Elsevier.
\item Fuller W.B., Thompson S.E. (1907). The laws of proportioning concrete. \emph{Trans. ASCE} 59, 67--143.
\item Kim Y.-R., Little D.N., Lytton R.L. (2003). Fatigue and healing characterization of asphalt mixtures. \emph{J. Mater. Civ. Eng.} 15(1), 75--83 (mortier bitumineux, FAM).
\item Westergaard H.M. (1939). Bearing pressures and cracks. \emph{J. Applied Mechanics} 6, A49--A53.
\item David L. (2026). Rapport de TFE, LTDS--ENTPE ; notice \cs{} (v0.2).
\end{itemize}
\normalsize\noindent\emph{Références à vérifier dans Zotero.}
\end{document}
